% concatenated from main.tex
\documentclass[letterpaper,reqno,12pt]{amsart} 
\usepackage{setspace}
\usepackage{graphicx} % Required for inserting images
\usepackage{cite}
\usepackage{amssymb, enumerate, color}
\usepackage{hyperref}
%\usepackage[dvipdfmx]{hyperref}
%\usepackage{pxjahyper}
\usepackage{enumitem}
\usepackage{graphicx}
\usepackage{subcaption}
\usepackage{here}
\usepackage[margin=1.2in]{geometry}
\usepackage{amsmath}
\usepackage{fancyhdr}
\usepackage{indentfirst}
\usepackage{tikz-cd}
\usepackage{tikz}
\usepackage{amsthm}
\usepackage[section]{placeins}
\usepackage{mathrsfs}
\usepackage[all]{xy}
\usepackage{setspace}
\usepackage{xcolor}
\usepackage[T1]{fontenc}
\usepackage[utf8]{inputenc}
\usepackage{cleveref}
\hypersetup{colorlinks=true,colorlinks=true,linkcolor=blue,urlcolor=blue, citecolor=blue}



\DeclareMathOperator{\Bs}{Bs}
\DeclareMathOperator{\Supp}{Supp}
\DeclareMathOperator{\Sym}{Sym}
\DeclareMathOperator{\CDiv}{CDiv}
\DeclareMathOperator{\CS}{CS}
\DeclareMathOperator{\reg}{reg}
\DeclareMathOperator{\Fix}{Fix}
\DeclareMathOperator{\LCS}{LCS}
\DeclareMathOperator{\bMov}{Mov}
\DeclareMathOperator{\NEbar}{\overline{\mathrm{NE}}}
\DeclareMathOperator{\NonSing}{NonSing}
\DeclareMathOperator{\cent}{center}
\DeclareMathOperator{\cont}{cont}
\DeclareMathOperator{\codim}{codim}
\DeclareMathOperator{\dival}{\mbox{$\mathcal R$}}
\DeclareMathOperator{\divi}{div}
\DeclareMathOperator{\flpal}{\mbox{$\mathcal{FR}$}}
\DeclareMathOperator{\im}{im}
\DeclareMathOperator{\mld}{mld}
\DeclareMathOperator{\mult}{mult}
\DeclareMathOperator{\Proj}{Proj}
\DeclareMathOperator{\Spec}{Spec}
\DeclareMathOperator{\exc}{exc}
\DeclareMathOperator{\ind}{index}
\DeclareMathOperator{\regular}{reg}
\DeclareMathOperator{\Accum}{Accum}
\DeclareMathOperator{\ima}{im}
\DeclareMathOperator{\kernel}{ker}
\DeclareMathOperator{\dom}{dom}
\DeclareMathOperator{\nil}{nil}
\DeclareMathOperator{\internal}{int}
\DeclareMathOperator{\Ass}{Ass}
\DeclareMathOperator{\Hom}{Hom}
\DeclareMathOperator{\Pic}{Pic}
\DeclareMathOperator{\lct}{lct}
\DeclareMathOperator{\Diff}{Diff}
\DeclareMathOperator{\chara}{char}

\newcommand{\red}{\textcolor{red}}
\newcommand{\tbf}{\textbf}
\newcommand{\tit}{\textit}
\numberwithin{equation}{section} 

\theoremstyle{plain}

\newtheorem{theorem}{Theorem}[section]
\newtheorem{lemma}[theorem]{Lemma}
\newtheorem{proposition}[theorem]{Proposition}
\newtheorem{theoremdef}[theorem]{Theorem-Definition}
\newtheorem{corollary}[theorem]{Corollary}
\newtheorem{conjecture}[theorem]{Conjecture}
\newtheorem{claim}[theorem]{Claim}
\newtheorem{definition}{Definition}[section]
\newtheorem{remark}{Remark}[section]
\newtheorem{example}{Example}[section]
%%%%%%%%%%%%%%%%%%%%%%%%%%%%%%%%%%%%%%%%%%%%%%
\title{Boundedness of complements for fibered Fano threefolds in positive characteristic}
%%%%%%%%%%%%%%%%%%%%%%%%%%%%%%%%%%%%%%%%%%%%%%
\author{Xintong Jiang}
\address{Tsinghua University}
\email{xt-jiang21@mails.tsinghua.edu.cn}
\date{}

\begin{document}
\begin{abstract}
    In this paper, we prove the canonical bundle formula for Fano type fibrations and Shokurov's conjecture on boundedness of complements for Fano type threefold pairs $(X,B)$ with fibration structures in large characteristics. In particular, we prove the conjecture when $-(K_X+B)\not\equiv 0$ is nef and not big in large characteristics.
\end{abstract}
\maketitle
\markboth{Xintong Jiang}{Boundedness of complements for fibered Fano threefolds in positive characteristic}
\tableofcontents
\section{Introduction}














%%%%%%%%%%%%%%%%%%%%%%%%%%%%%%%%%%%%%%%%%%%%%%%%%%%%%%%%%%%%%%%%%%%%%%%%%%%%%%%%%%%%%
We work over an algebraically closed field $k$ of characteristic $p>5$ through this paper. In birational geometry, a central question is whether a  given class of varieties is bounded. Proving that a class $\mathcal P$ of varieties is bounded allows for a more explicit construction and study of its moduli space, yielding results that facilitate further classification. In characteristic zero, many important boundedness results are known. The boundedness of varieties of general type is proved in \cite{hacon2012birationalautomorphismsvarietiesgeneral}\cite{HMXboundednessofgeneraltype}, the boundedness of Fano varieties with mild singularities is established by a sequence of work \cite{birkar2019antipluricanonicalsystemsfanovarieties}\cite{birkar2020singularitieslinearsystemsboundedness}\cite{birkarfanofib}. For Calabi-Yau varieties, the problem is hard since the lack of a canonical polarization. However, there are still some satisfactory results in this direction \cite{birkarpolarisedvarieties}\cite{JJZboundedpolarizedLCYfib}. In positive characteristic, very few results on boundedness are known even for 3-dimensional varieties.  The difficulty essentially comes from the failure of vanishing theorems and the lack of understanding of very special morphisms that appear only in positive characteristic like inseparable morphisms. However, for Fano varieties in positive characteristic, especially in large characteristics, many pathologies might be controlled which enables us to study them more easily. In particular, the following famous BAB conjecture is widely supposed to be true in large characteristics:
\begin{conjecture}[BAB Conjecture]
\label{BAB}
    Let $d$ be a natural number and $\epsilon>0$ be a real number, then the set of $\epsilon$-lc Fano varieties of dimension $d$ forms a bounded family.
\end{conjecture}
The condition $\epsilon$-lc cannot be strengthened into klt, cf. \cite[1.2]{birkar2020singularitieslinearsystemsboundedness}. This conjecture was proved in characteristic 0 by Birkar in   \cite{birkar2019antipluricanonicalsystemsfanovarieties}\cite{birkar2020singularitieslinearsystemsboundedness} using Shokurov's theory of complements. In positive characteristic, the conjecture was proved in dimension 2 by Alexeev \cite{boundednessandk2} and for toric varieties by the Borisov brothers \cite{Borisov1993SINGULARTF}. In dimension 3, the conjecture is  widely open with very few limited results, see  \cite{sato2024boundednessweakfanothreefolds}\cite{Das_2019} and etc.

According to the methods in  \cite{birkar2019antipluricanonicalsystemsfanovarieties}\cite{birkar2020singularitieslinearsystemsboundedness}, a proposition closely related to the BAB conjecture is the boundedness of complements on Fano type varieties:
\begin{conjecture}[Boundedness of complements, Shokurov]
\label{cplt}
    Let $d,p\in\mathbb N$ be natural numbers and $R\subset[0,1]$ be a finite set of rational numbers, then there exists a constant $n$ only depending on $d,p$ and $R$ such that if
    \begin{enumerate}
        \item $(X',B'+M')$ is a projective lc generalized pair of dimension $d$,
        \item $X\to Z$ is a contraction,
        \item $B'\in \Phi(R)$, $pM$ b-Cartier,
        \item $X'$ Fano type over $Z$, and
        \item $-(K_{X'}+B'+M')$ is nef over $Z$.
        
    \end{enumerate}Then $K_{X'}+B'+M'$ has an $n$-complement $K_{X'}+B^{\prime+} +M'$ with $B^{\prime+}\geq B'$ over any point $z\in Z$.
\end{conjecture} 

The conjecture is known for $\dim X=2$ in all characteristic since \Cref{BAB} is known in full generality \cite{boundednessandk2}. In characteristic $0$, the conjecture was proved by induction on dimensions, separated into two cases, the exceptional case and the non-exceptional case. In the non-exceptional case, one of the main step is to construct complements from a fibration structure using induction hypothesis. 

In this paper, we will prove the following statement which predicts that the same process would work in large characteristic.
\begin{theorem}[\Cref{fibcomp}]
\label{fibcomp1}
    Let $R\subset [0,1]$ be a finite set of rational numbers, there is an $n\in\mathbb N$ and a prime number $p_0$ depending only on $R$ such that if $(X,B)$ is a projective lc pair defined over an algebraically closed field $k$ satisfying the following conditions:
    \begin{enumerate}
        \item $\chara k=p>p_0$,
        \item $-(K_X+B)$ is nef,
        \item $X$ is of Fano type,
        \item $B\in \Phi(R)$ and
        \item there is a contraction $f:(X,B)\to V$ such that $(K_X+B)\sim_{\mathbb Q} 0/V$ with $3>\dim V>0$,
    \end{enumerate}
    then there is an $n$-complement $(X,B^+)$ for $(X,B)$ with $B^+\geq B$.
\end{theorem}
A direct corollary is the following theorem, the result hopefully holds for generalized pairs.
\begin{theorem}[\Cref{fibcomp2}]
\label{Main}
       

    Let $R\subset [0,1]$ be a finite set of rational numbers. Then there is a constant $n$ and a prime number $p_0$ depending only on $R$, such that if $(X,B)$ is a $3$-dimensional projective lc pair defined over an algebraically closed field $k$ satisfying the following conditions:
    \begin{enumerate}
        \item $\chara k=p>p_0$,
        \item $B\in \Phi(R)$,
        \item $X$ is Fano type and
        \item $-(K_X+B)\not\equiv0$ is nef but not big,
    \end{enumerate}
    then there is an $n$-complement $(X,B^+)$ with $B^+\geq B$.
\end{theorem}

The strategy for proving this theorem is to use the canonical bundle formula and induction. Though the canonical bundle formula generally does not hold in positive characteristic, we can still obtain some results for Fano type fibrations:
\begin{theorem}[Canonical bundle formula for Fano type fibrations, \Cref{cbf1}]
\label{cbf}
    Assume $(X,B)$ is a lc pair over an algebraically closed field of characteristic $p$, $f:X\to Z$ is a contraction with dim $X=3$ and dim $Z>0$, $K_X+B\sim_{\mathbb Q} 0/Z$, $B$ is relatively big over $Z$. Let $\epsilon>0$ be a real number and $R\subset [0,1]$ be a finite set of rational numbers, then there is a prime number $p_0=p_0(\epsilon,R)$ such that if $\chara k=p>p_0$ and one of the following conditions holds:
    \begin{enumerate}
        \item $\dim Z=3$, i.e. $X\to Z$ is birational,
        \item $\dim Z=2$, $B^{h}\in\Phi(R)$ where $B^h$ is the horizontal part of $B$,
        \item $\dim Z=2$, $(X,B)/Z$ has lc general fibers,
        \item $\dim Z=1$, $(X_\eta,B_\eta)$ is $\epsilon$-lc and $B^{h}\in\Phi(R)$,
        \item $\dim Z=1$, $X_\eta$ is an $\epsilon$-lc Fano variety and $B^h\in\Phi(R)$.
    \end{enumerate}
    Then we have the following formula:
    \[K_X+B\sim_{\mathbb Q} f^*(K_Z+B_Z+M_Z)\]
    where \[B_Z:=\sum_{D\text{ prime divisor on }Z}(1-\text{lct}_{\eta(D)}(f^*D,X,B))D\]
    is the discriminant part and $M_Z$ is a b-nef ${\mathbb Q}$-b-divisor.
\end{theorem}

The canonical bundle formula connects the singularities and positivity of the base space and the whole space of a fibration, in particular, we will have the following pleasant corollary:
\begin{corollary}[Contraction of Fano type threefolds is of Fano type, \Cref{contr}]
\label{contrac}
    Assume $(X,B)$ be a $3$-dimensional projective lc pair with a contraction $f:X\to Z$, where $K_X+B\sim_{{\mathbb Q}}0/Z$. Let $R\subset [0,1]$ be a finite set of rational numbers and $\epsilon>0$ be a real number, then there exists a prime number $p_0$ depending only on $R$ and $\epsilon$ such that if $-(K_X+B)$ is nef, $X$ is of Fano type, char $k=p>p_0$, $B\in\Phi(R)$ and one of the following conditions holds:
    \begin{enumerate}
        \item $Z$ is a projective normal surface,
        \item $Z$ is a projective normal curve, $(X,B)$ is $\mathbb Q$-factorial dlt and $X_\eta$ is an $\epsilon$-lc Fano surface.
    \end{enumerate}
    Then $Z$ is of Fano type, except in the case where $X\to Z$ is a surface fibration over an elliptic curve and the canonical bundle formula takes the form 
    \[K_X+B\sim_\mathbb Qf^*K_Z\sim_\mathbb Q0.\]
    Here $B_Z=M_Z=0$, $(X,B)$ is not klt and not geometrically klt, $B^\text{ver}=0$ and each normalization $S^\nu$ of a lc place $S\in\lfloor B\rfloor$ is not of Fano type.
\end{corollary}

This corollary is essential for the induction steps in the proof of \Cref{Main}. We also prove that the following results which predicts that relative complements for Fano type fibrations in large characteristic is bounded, which is a crucial input for the boundedness of global complements.

\begin{theorem}[\Cref{recomp}]
\label{relcomp}
    Assume $(X,B)$ is a 3-dimensional projective lc pair, $f:X\to Z$ is a contraction with $3>\dim Z>0$, let $R\subset [0,1]$ be a finite set of rational numbers, then there is some natural number $n=n(R)\in\mathbb N$ a prime number $p_0=p_0(R)$ such that suppose:
    \begin{enumerate}
        \item $\chara k>p_0$
        \item $B\in\Phi(R)$,
        \item $X$ is of Fano type$/Z$,
        \item $-(K_X+B)$ is nef$/Z$.
    \end{enumerate}
    Then, for any $z\in Z$, there is an $n$-complement $K_X+B^+$ of $K_X+B$ over $z$ with $B^+\geq B$.
\end{theorem}



We now give an outline of the paper. In Chapter 2, we will give the preliminary knowledge on birational geometry. In Chapter 3, we prove \Cref{cbf} with a similar method in \cite{jxtcbf} by proving generic normality. In Chapter 4, we prove one side of our main theorems \Cref{fibcomp1} and \Cref{Main} by using some partial result of \Cref{relcomp}, the main techniques are a limited vanishing theorem for Fano type fibrations in positive characteristic and the effective canonical bundle formula. In Chapter 5, we finish the proof of \Cref{relcomp} by lifting Frobenius stable sections. We will see the main difficulty of \Cref{relcomp} comes from the case when $-(K_X+B)$ is nef and big over a curve $Z$. If $Z$ is replaced by a point, this is the question on the boundedness of global complements for lc weak Fano pairs, which has not yet been solved.

Many of the ideas in this article are modified and retrofitted from \cite{birkar2019antipluricanonicalsystemsfanovarieties} and some other papers on birational geometry in positive characteristic. If some work’s proof is absolutely characteristic free but the known reference only proves the 0-characteristic case, we will refer it directly.

\paragraph{\tbf{Acknowledgments}}
    The author would like to thank his supervisor Caucher Birkar for his patient guidance and suggestions. The author also would like to thank Zheng Xu, Fulin Xu, Xiaowei Jiang, Haoran Zhu, Florin Ambro for reading the draft of this paper and giving the author many important comments while writing this paper.





\section{Preliminaries}









%%%%%%%%%%%%%%%%%%%%%%%%%%%%%%%%%%%%%%%%%%%%%%%%%%%%%%%%%%%%%%%%%%%%%%%%%%%%%%
In this chapter, some basic notions and results of birational geometry are mentioned for readers who are not familiar with the subject. In this paper, all varieties are quasi-projective and reduced schemes of finite type over an algebraically closed field $k$ of characteristic $p>5$ and the ambient variety $X$ is projective unless stated otherwise. 

\paragraph{\tbf{Resolution of singularities}}
For a variety $X$, a resolution of singularity is a proper birational map $f:Y\to X$ from a smooth variety $Y$, which is an isomorphism over the regular locus of $X$, and for the singular locus $X_{sing}$, one has $f^{-1}X_{sing}$ is a divisor with simple normal crossings. 
\begin{theorem}
    For a 3-dimensional variety $X$, there is a resolution of singularity $f:Y\to X$ which is obtained by a sequence of blow-ups along smooth centers over $X_{sing}$.
\end{theorem}
\begin{proof}
    See \cite{Cutkosky2004ResolutionOS}, \cite{Cossart2008ResolutionOS} and \cite{Cossart2009RESOLUTIONOS}.
\end{proof} 

\paragraph{\tbf{Hyperstandard sets}}
Let $R\subset [0,1]$ be a finite set of rational numbers, we define the hyperstandard set $\Phi(R)$ associated to $R$ to be:
\[\Phi(R):=\{0,1\}\cup\{1-\frac{r}{m}|r\in R, m\in \mathbb Z^+\}.\]
$\Phi(R)$ is a DCC set (i.e., a set of real numbers in which every non-empty subset has a minimal element) of rational numbers contained in $[0,1]$ with the only accumulation point $1$.

If $\Phi(R)$ is a hyperstandard set and $I$ is a common denominator of elements in $R$, then for any $b\in \Phi(R)$ and $n$ divisible by $I$, $nb\leq\lfloor (n+1)b\rfloor$.

\paragraph{\tbf{Divisors, pairs and linear systems}} We define $\mathbb N=\mathbb Z^{\geq0}$. Divisors refer to Weil divisors, i.e., formal linear combinations of integral codimension-1 closed subvarieties. These correspond to reflexive sheaves of rank one up to linear equivalence. The notion of Cartier divisors follows the usual definition, which correspond to invertible sheaves or line bundles up to linear equivalence. A divisor is called a prime divisor if it is integral. A divisor $D$ over $X$ is defined as a divisor on some birational model $W\to X$. For a variety $X$, the dualizing sheaf $\omega_X$ denotes to be the lowest cohomology sheaf of its dualizing complex. The reflexive hull of $\omega_X$ corresponds to a divisor $K_X$ up to linear equivalence, which is called the canonical divisor of $X$. 

For $\mathbb F=\mathbb Q,\mathbb Z_{(p)}=(\mathbb Z-p\mathbb Z)^{-1}\mathbb Z$, $\mathbb F$-divisors are the $\mathbb F$-linear combination of divisors and $\mathbb F$-linear equivalences between divisors are generated by $\mathbb F$-linear combination of linear equivalences. A similar definition applies for $\mathbb F$-Cartier divisors. An $\mathbb F$-divisor is $\mathbb F$-Cartier if it is $\mathbb F$-linear equivalent to an $\mathbb F$-Cartier divisor. $(X,B)$ is called a sub pair if $X$ is a normal variety and $B$ is a $\mathbb Q$-divisor such that $K_X+B$ is $\mathbb Q$-Cartier and $B\leq 1$ (in coefficients).  A sub pair $(X,B)$ is called a pair if $B\geq 0$. Here, $B$ is called a boundary if $(X,B)$ is a pair.

For an $\mathbb F$-divisor $M$, we often denote \[H^0(M):=H^0(X,\mathcal O_X(\lfloor M\rfloor))=\{f\in k(X)\mid\text{div}(f)+M\geq 0\}.\] The linear system of $M$ is defined as  \[|M|:=\Proj(H^0(M))=\{N\sim M\mid N\geq0\}.\] The $\mathbb F$-linear system is defined as \[|M|_{\mathbb F}:=\{N\sim_{\mathbb F}\mid N\geq0\},\]in particular \[|M|_{\mathbb Q}=\bigcup\limits_{m\in\mathbb N} \frac{1}{m}|mM|.\]The base locus $Bs(|M|)$ denotes the maximal closed subset of $X$ contained in each $N\in|M|$, the stable base locus is defined as \[Bs(|M|_{\mathbb Q}):=\bigcap\limits_{m\in\mathbb N} Bs(|mM|).\] If $|M|\neq\emptyset$, $|M|$ will define a rational map \[\phi_M:X\dashrightarrow |M|^\lor\simeq \mathbb P^n,\] which is determined on $X\setminus Bs(|M|)$ by mapping $x$ to the hyperplane in $H^0(M)$ consisting of sections whose corresponding divisors $N\sim M$ pass through $x$. 

\paragraph{\tbf{B-divisors and generalized pairs}}

Let $X$ be a variety, a b-divisor $\mathbf M$ on $X$ is a configuration of divisors $\mathbf M_Y$ on each projective birational model $Y$ over $X$ such that if $f:Z\to Y$ is a morphism of birational models over $X$, then $f_*\mathbf M_Z=\mathbf M_Y$. A b-divisor is said to be represented by $Y$ if $\mathbf M_Y$ is $\mathbb R$-Cartier and if $\mathbf M_Z=f^*\mathbf M_Y$ holds for any projective birational morphism $f:Z\to Y$. Usually we will use $M:=\mathbf M_X$ to represent the b-divisor $\mathbf M$ for convenience.

We say a b-divisor $\mathbf M$ is b-nef if it is represented by some model $Y$ and $\mathbf M_Y$ is nef. Suppose $X$ is a $\mathbb Q$-factorial surface and $\mathbf M$ is a b-nef b-$\mathbb Q$-divisor on $X$, then $M$ is nef. Indeed, if $\mathbf M$ is represented by some $Y$ and $C$ is any curve on $X$, then \[M.C=f_*\mathbf M_Y.C=\mathbf M_Y.f^*C\geq0,\] which implies $M$ is nef. In general, a b-nef b-divisor $\mathbf M$ is not nef on $X$.

A generalized pair is given as $(X',B'+M')/Z$ where $X'$ is a normal variety with a projective morphism $X'\to Z$, $B'\geq0$ a $\mathbb Q$-divisor (usually $B'\leq 1$) on $X'$ and a b-$\mathbb Q$-Cartier $b$-$\mathbb Q$-divisor $M'$ represented by some projective birational morphism $\phi:X\to X'$ and a $\mathbb Q$-Cartier $\mathbb Q$-divisor $M$ on $X$ such that $M$ is nef over $Z$ and $M'=\phi_*M$ and $K_{X'}+B'+M'$ is $\mathbb Q$-Cartier. Since $M'$ is defined birationally, one may assume that $X\to X'$ is a log resolution. $M$ is viewed as a b-divisor in generalized pairs.

\paragraph{\tbf{Singularities from MMP}}
Suppose $D$ is a prime divisor on $X$, for any $\mathbb Q$-divisor $A$ on $X$ we define $\mu_D(A)$ to be the coefficient of $D$ in $A$. For a prime divisor $D$ on a log resolution $W/X$ of the (resp. sub-)pair $(X,B)$, let $K_W+B_W$ be the pullback of $K_X+B$, the log discrepancy of $(X,B)$ is defined as \[a(D,X,B):=1-\mu_D(B_W).\]The log discrepancy is a number defined up to strict transformations along the birational maps between smooth models of $X$. One say the (resp. sub-)pair $(X,B)$ is (resp.sub-)lc (resp. klt, plt, canonical, terminal, $\epsilon$-lc) if $a(D,X,B)\geq 0$ (resp. $>0$, $>0$ for exceptional $D$, $\geq 1$ for exceptional $D$, $>1$ for exceptional $D$, $\geq\epsilon$) for every $D$ over $X$. A non-klt place of a sub pair $(X,B)$ is a prime divisor $D$ on birational models of $X$ such that $a(D,X,B)\leq0$. A non-klt center is the image on $X$ of a non-klt place. A (resp.sub-)pair is (resp.sub-)dlt if it is lc and log smooth near generic points of non-klt centers.

For a generalized pair $(X',B'+M')$ and a divisor $D$ over $X$, take a sufficiently high resolution $f:X\to X'$ defining $M'=f_*M$ and contains $D$, we define $K_X+B+M:=f^*(K_{X'}+B'+M')$, and one can similarly define generalized version of lc, klt, plt, $\epsilon$-lc by considering the generalized log discrepancy \[a(D,X',B'+M'):=1-\mu_D(B).\]If $M=0$, these notions of singularities will coincide with the classical version.

\paragraph{\tbf{Contractions and minimal model programs}}
An algebraic fiber space, or a contraction, $f:X\to Y$ is a projective morphism  between varieties satisfying $f_* \mathcal O_X=\mathcal O_Y$. This condition is equivalent to $f$ being a projective surjective morphism whose Stein factorization has a trivial finite part, or alternatively, to $f$ being a projective surjective morphism such that the function field of $Y$ is algebraically closed in that of $X$. It's well known the fibers are connected.

We use standard results of the minimal model program (MMP), MMP in char $k>5$ up to dimension 3 is already fully known in \cite[1.1]{KYH20lcmmp}:
\begin{theorem}[cf. \cite{KYH20lcmmp}]
    Let $(X,B)/Z$ be a 3-dimensional lc pair over $k$ of char $>5$, $X\to Z$ be a projective contraction, then there is a minimal model program$/Z$ on $K_X+B$ such that:
    \begin{enumerate}
        \item If $K_X+B$ is pseudo-effective$/Z$, then the MMP terminates with a log minimal model$/Z$.
        \item If $K_X+B$ is not pseudo-effective$/Z$,then the MMP terminates with a Mori fiber space$/Z$.
    \end{enumerate}
\end{theorem}

\paragraph{\tbf{$\mathbb Q$-factorialization}}
A normal variety is $\mathbb Q$-factorial if every divisor is $\mathbb Q$-Cartier. For a generalized pair $(X',B',M')$ with data $X\to X'$, a $\mathbb Q$-factorial generalized dlt model is a $\mathbb Q$-factorial generalized dlt generalized pair $(X'',B''+M'')$ with a projective birational morphism $\psi:X''\to X'$ under a log resolution $X\to X''$ (after taking a common resolution) such that $B''$ and $M''$ are pushdowns of $B$ and $M$, in particular $K_{X''}+B''+M''=\psi^*(K_{X'}+B'+M')$, and if every exceptional prime divisor of $\psi$ appears in $B''$ with coefficient 1. Such model exists for generalized lc pairs. If $(X',B'+M')$ is generalized klt, then there is a $\mathbb Q$-factorial generalized klt model and $\psi$ is a small morphism (i.e. no divisor is contracted or extracted), this is called a small $\mathbb Q$-factorialization.

\paragraph{\tbf{Volumes, Kodaira dimensions and Iitaka fibrations}}
Let $X$ be a normal projective variety of dimension $d$, and let $D$ be a $\mathbb Q$-divisor on $X$. The Kodaira dimension $\kappa(D)$ (resp. the numerical Kodaira dimension $\kappa_\sigma(D)$) is defined as $-\infty$ if $D$ is not effective (resp. pseudo-effective), and otherwise as the largest integer $r$ such that \[\limsup_{m\to\infty}\frac{h^0(\lfloor mD\rfloor)}{m^r}>0.\] resp. for some very ample divisor $A$\[\limsup_{m\to\infty}\frac{h^0(\lfloor mD\rfloor+A)}{m^r}>0.\] We define the volume \[\text{vol}(D):=\limsup_{m\to \infty}\frac{h^0(\lfloor mD\rfloor)}{m^d},\] and say $D$ is big if $vol(D)>0$ as usual. $|\lfloor mD\rfloor|$ will define a morphism $\phi_m=\phi_{\lfloor mD\rfloor}$. The dimension of the image of $\phi_m$ will stabilize to $\kappa(D)$. The stabilized rational fibration $\phi:X\dashrightarrow \phi(X)$ is called the Iitaka fibration.

\paragraph{\tbf{Fano pairs and varieties of Fano type}}
Let $(X,B)$ be a pair with a contraction $X\to Z$, we say $(X,B)$ is log Fano (resp. weak log Fano) over $Z$ if $-(K_X+B)$ is ample (resp. nef and big) over $Z$. We assume $B=0$ when we don't mention $B$. We say a pair $(X,\Delta)$ is of Fano type over $Z$ if there exists a boundary $B$ such that $(X,\Delta+B)$ is a klt weak log Fano pair over $Z$, or equivalently, if there exists a boundary $\Gamma$ such that $(X,\Delta+\Gamma)$ is klt, $\Gamma$ is big over $Z$, and $K_{X}+\Delta+\Gamma\sim_{\mathbb{Q}}0/Z$.

Suppose $f:X\to Y$ is a birational contraction and $(X,B)$ is of Fano type, then $(Y,B_Y=f_*B)$ is of Fano type as the pushforward of a big divisor is big. Suppose $(X,B)\dashrightarrow(Y,B_Y)$ is a sub-crepant birational map (i.e. there is a common resolution $W$ such that $(K_X+B)|_W\leq (K_Y+B_Y)|_W$), then $(Y,B_Y)$ is of Fano type will imply that $(X,B)$ is of Fano type. Hence taking crepant resolutions, running MMP and taking $\mathbb Q$-factorializations will keep the property of Fano type.

Let $X$ be a variety $\mathbb Q$-factorial of Fano type and suppose $(X,\Delta)$ is klt and $K_X+\Delta\sim_{\mathbb Q}0$. Let $D$ be a $\mathbb Q$-divisor on $X$, then for $\epsilon\ll 1$, we have \[\epsilon D\sim_{\mathbb Q} K_X+\Delta+\epsilon D\sim_{\mathbb Q}K_X+(1-n\epsilon)\Delta +n\epsilon \Delta+\epsilon D.\] Since $\Delta$ is big, we have $\Delta\sim_{\mathbb Q} B+A$ for some effective $\mathbb Q$-divisor $B$ and ample $\mathbb Q$-divisor $A$ on $X$, then we can always find some $n$ large enough such that there is some $ H\sim_{\mathbb Q} nA$ and $H+D>0$. So 
\[\epsilon D\sim_{\mathbb Q}K_X+(1-n\epsilon)\Delta +n\epsilon \Delta+\epsilon D\sim_{\mathbb Q}K_X+(1-n\epsilon)\Delta+ n\epsilon B+\epsilon (H+D).\]
Since $(X,\Delta)$ is klt, for $\epsilon\ll 1$, we always have $(X,(1-n\epsilon)\Delta+ n\epsilon B+\epsilon (H+D))$ is klt. Hence we can always run $\epsilon D$-MMP on $X$ which will terminate to get some models. Moreover if $D$ is a nef divisor on a threefold of Fano type of $\chara>5$, then $D$ is semi-ample by base-point free theorem \cite[1.2]{birkar2014existencemorifibrespaces}.

\paragraph{\tbf{Bounded families}}
Now we introduce the notion of bounded families mentioned in the BAB conjecture. A couple $(X,D)$ is formed by a normal projective variety $X$ and a divisor $D$ on $X$ such that the coefficient of $D$ falls in $\{0,1\}$. Isomorphisms between couples are isomorphisms between base schemes such that the morphism is compatible and onto for boundaries. 

A set $\mathcal P$ of couples is birationally bounded (resp. bounded) over a scheme $S$ if there exist finitely many projective flat morphisms $V^i\to T^i$ of integral schemes of finite type over $S$ and reduced divisors $C^i$ on $V^i$ such that for each $(X,D)\in\mathcal P$ there is an $i$ and a closed point $t=\Spec(H^0(X,\mathcal O_X))\in T^i$ and a birational isomorphism (resp. isomorphism) $\phi:V_t^i\dashrightarrow X$ such that the fiber $(V_t^i,C_t^i)$ over $t$ is a couple and $E\leq C_t^i$, where $E$ is the sum of the strict transform of $D$ and the reduced exceptional divisor of $\phi$. A set $\mathcal R$ of projective pairs $(X,B)$ is said to be log birationally bounded (resp. log bounded) over $S$ if the set of $(X,\text{Supp}B)$ is birationally bounded (resp. bounded) over $S$. And if $B=0$ for all the elements in $\mathcal R$, we usually remove the log and say the set is birationally bounded (resp. bounded) over $S$. If $S=\Spec(k)$ is the base field we are working on, usually we omit the suffix over $S$ and simply say $\mathcal P$ is birationally bounded. We offer a useful characteristic-free criterion for boundedness here, and from this one can assume that $(V_t^i,C_t^i)$ is isomorphic to $(X,D)$ for bounded families in the above definition (cf. \cite[2.21]{birkar2019antipluricanonicalsystemsfanovarieties}).
\begin{lemma}[\cite{birkar2019antipluricanonicalsystemsfanovarieties} 2.20]
    If $\mathcal P$ is a set of couples of dimension $d$, $\mathcal P$ is bounded if and only if there is an $r\in \mathbb N$ such that for any $(X,D)\in \mathcal P$, there is a very ample divisor $A$ on $X$ such that $A^d\leq r$ and $A^{r-1}D\leq r$.
\end{lemma}
If a set of varieties $\mathcal R$ is bounded, then the Gorenstein indices, the (anti-)canonical volumes, the indices of the effective Iitaka fibrations (e.g. for varieties that $\kappa(K_X)\geq 0$, the minimal $m$ such that $|mK_X|$ defines the Iitaka fibration), the Picard numbers, etc. are all bounded.

\paragraph{\tbf{Complements}}

Now we introduce the terminology of complements, which is introduced by Shokurov in \cite{3foldflip}. Let $(X',B'+M')/Z$ be a generalized pair, set $T':=\lfloor B'\rfloor$ and $\Delta':=B'-T'$. An $n$-complement of $K_{X'}+B'+M'$ over $z\in Z$ is of the form $K_{X'}+B'^++M'$ such that over some neighbourhood of $z$, $(X,B'^++M')$ is generalized lc, $nM$ is b-Cartier, and \[n(K_{X'}+B'^++M')\sim 0,nB'^+\geq nT'+\lfloor (n+1)\Delta'\rfloor.\] Moreover if $(X',B'+M')$ is generalized klt, then we say $B^+$ is a klt $n$-complement. The boundedness of complements is highly related to the boundedness of varieties. Indeed, the BAB conjecture in characteristic $0$ is proved by proving that these Fano type varieties admits bounded klt complements \cite{birkar2020singularitieslinearsystemsboundedness}.

\paragraph{\tbf{Vanishing theorem for pl-contraction of plt pairs}}
The Kawamata vanishing theorem is generally not true in positive characteristic even for surfaces, a famous counterexample is the Raynaud surfaces, one can also find a counterexample in \cite{Counterexamplevanishing}. However, for log del Pezzo surfaces and smooth Fano threefolds, there is a sequence of work which shows that some limited version of vanishing theorems works. Here we give a vanishing theorem for pl-contraction of plt pairs in positive characteristic which is useful to prove the boundedness of relative complements:
\begin{lemma}\label{van}
    Assume $f:(X,\Gamma)\to Z$ is a contraction of projective varieties, $\phi:X'\to (X,\Gamma)$ is a log resolution, $D$ is a Weil divisor on $X'$, $\Lambda'$ is an effective $\mathbb Q$-divisor on $X$ such that:
    \begin{enumerate}
        \item dim $X=3$, dim $Z>0$, char $k>5$
        \item $(X,\Gamma)$ and $(X',\Lambda')$ are $\mathbb Q$-factorial plt with an irreducible lc center $S:=\lfloor\Gamma\rfloor$ and $S':=\lfloor \Lambda'\rfloor$, $S'$ is the strict transformation of $S$,
        \item $S$ is vertical$/Z$,
        \item $-(K_X+\Gamma)$ is ample$/Z$,
        \item $-S$ is nef$/Z$ if $Z$ is a curve,
        \item $D\sim_{\mathbb R}K_{X'}+\Lambda'-S'+L'$,
        \item $K_X+\Lambda:=\phi_*(K_{X'}+\Lambda')\leq K_X+\Gamma$,
        \item $L'=\phi^*L$ and $L$ is an ample $\mathbb Q$-divisor$/Z$.
    \end{enumerate}
    Then $R^1f'_*\mathcal O_X(D)=0$ near $f(S)$.
\end{lemma}
\begin{proof}
    By Grauert-Riemenschneider vanishing for 3-dimensional dlt excellent pairs with char $>5$ \cite[3]{BernasconiKollarGRvanishing}, $R^i\phi_*(D)=0$ for $i>0$. If $-S$ is nef$/Z$, we see \[R^1f'_*(D)=R^1f_*(\phi_*(D))=R^1f_*(K_X+\Lambda+(L-S))=0\text{ near }f(S)\] by \cite[1.1]{Bernasconiplvanishing} and the Leray spectral sequence. If $Z$ is a surface, then proofs of \cite[2.2]{Bernasconiplvanishing} and \cite[2.8]{Bernasconibasepointfreeklt} apply similarly for $K_X+\Lambda+L-S$, hence again \[R^1f'_*(D)=R^1f_*(K_X+\Lambda+(L-S))=0\text{ near }f(S)\].
\end{proof}

\paragraph{\tbf{Singularities near the generic fiber}}
Suppose $X\to Z$ is a fibration, then any $\mathbb Q$-divisor $B$ can be uniquely factored as the sum of two $\mathbb Q$-divisors:
\[B=B^h+B^v.\]
Here $B^h$ is called the horizontal part, which contains all the components of $B$ which is dominant over $Z$, and $B^v$ is called the vertical part, consisting of the components of $B$ which is not dominant over $Z$. Let $\eta:=\eta_Z$ denotes the generic point of $Z$, clearly vertical part has no affects to the generic fiber:
\[(X_\eta,B_\eta):=(X,B)\times_Z\eta.\]
Moreover, the coefficients of $B_\eta$ is the same as the coefficients of $B^h$. One can define the generic log discrepancies for any divisor $D$ over $X_\eta$ as an $\eta$-variety to be
\[a_\eta(D,X,B):=a(D,X_\eta,B_\eta).\]
We define the generic version of singularities (klt, lc, etc) with the generic log discrepancies similarly. For Fano type contractions, we have the following lemma which shows small generic log discrepancies must be $0$:
\begin{lemma}
\label{singgen}
    Let $\Phi\subset [0,1]$ be a DCC set and $p>5$, then there is $\epsilon>0$ depending only on $\Phi$ such that if the contraction $(X,B)/Z$ is a 3-dimensional projective pair over a normal curve and $D$ is a prime divisor over $X_\eta$ satisfying that
    \begin{enumerate}
        \item $(X,B)$ is dlt near the generic fiber and $(X,0)$ is klt of Fano type near the generic fiber,
        \item $K_X+B\sim_{\mathbb Q}0/Z$ and $B^h\in\Phi$
        \item $a_\eta(D,X,B)<\epsilon$,
    \end{enumerate}
    then $a_\eta(D,X,B)=0$.
\end{lemma}
\begin{proof}
    We prove this by contradiction. Suppose $(X_i/Z_i,B_i,D_i)$ is a sequence of pairs and divisors such that $a_\eta(D_i,X_i,B_i)=\epsilon_i\to 0$. We set $X'_{i,\eta}\to X_{i,\eta}$ to be the morphism only extracts $D_i$ and let $K_{X'_{i,\eta}}+B_i'$ to be the pull back of $K_{X_{i,\eta}}+B_{i}$ and let $b_i:=1-\epsilon_i$. Since $X_i$ is of Fano type and $b_i>0$, it follows that $X'_i$ is also of Fano type and $B'_{i,\eta}$ is big. Since $B'_i\in\Phi':=\Phi\cup\{b_i|i\in\mathbb N\}$ is DCC and $\{b_i\}$ is not finite, we get a contradiction by global ACC for Fano type fibrations \cite[5.1]{waldron3mfsimperf}.
\end{proof}

\paragraph{\tbf{Generic normality of bounded families$/\mathbb Z$}}
The following facts are useful to prove the canonical bundle formula in positive characteristic.
\begin{lemma}\label{genorm}
    Suppose $\mathcal P$ is a bounded family$/\Spec\mathbb Z$ of projective $k$-varieties for $k$ varying among all fields. Suppose a component $X_0\subset X\in\mathcal P$ is a normal variety$/k$, then there is $p_0=p_0(\mathcal P)$ such that if char $k>p_0$, $X_0$ is geometrically normal.
\end{lemma}
\begin{proof}
    Since $\mathcal P$ is bounded, we see there is a flat proper morphism $\mathcal X\to\mathcal T$ of reduced separated schemes of finite type over $\Spec\mathbb Z$. Since $\mathcal T$ is excellent, we assume components of $\mathcal T$ are affine regular by taking a stratification and by Noetherian induction. $\mathcal T$ has finitely many irreducible components, consider $k\to \mathcal T$ parametrizing $X$, if it falls in an irreducible component $\mathcal S$ with the function field of characteristic $p>0$, we just ignore this component $\mathcal S$ and make $p_0>p$. So we may assume $\mathcal T$ is irreducible with function field of characteristic $0$. Now consider the normalization $\mathcal X^\nu$ of $\mathcal X$. We see that geometric normality, geometric regularity and geometric reducedness is an open condition (cf. \cite[II.6.9.1 and III.12.1.1]{PMIHES_1967__32__5_0}). Since $\mathcal X^\nu\to \mathcal T$ is a flat surjective morphism between integral schemes of characteristic $0$, outside a closed subscheme $\mathcal V\subset \mathcal T$, the fiber $X'$ of $\mathcal X^\nu/\mathcal T$ is geometrically normal and isomorphic to the normalization of the corresponding fiber of $\mathcal X/\mathcal T$.
        
    If $X_0$ is contained in the fiber over $t\in\mathcal T-\mathcal V$ of $\mathcal X/\mathcal T$. Since $X_0$ is normal and we have a morphism $X_0\to\mathcal X_t$ which is dominant over a component of $\mathcal X_t$, we have a natural morphism $f:X_0\to(\mathcal X^{\nu})_t$ by the universal property of a normalization. We claim that the morphism $X_0\to(\mathcal X^\nu)_t$ identifies $X_0$ as one component of the fiber $X':=(\mathcal X^\nu)_t$ over an open subset of $\mathcal T$. Indeed since the generic fiber of $\mathcal X/\mathcal T$ and $\mathcal X^\nu/\mathcal T$ is geometrically reduced, so is a geometric fiber over an open subset of $\mathcal T$, say outside $\mathcal W$. Since $\mathcal X/\mathcal T$ and $\mathcal X^\nu/\mathcal T$ are proper flat, we may assume dim $X$ is the relative dimension of $\mathcal X/\mathcal T$ and the families is equidimensional over an open subset of $\mathcal T$, say outside $\mathcal Y$. We now consider the case when $k\in\mathcal T-\mathcal W-\mathcal Y-\mathcal V$, we see $X_0\to X'\to X$ is identity on $X_0$. Since they have the same dimension, this gives a birational morphism between $X_0$ and a component of $X'$. Since $X_0$ is normal and $X'\to X$ is finite since $\mathcal X^\nu\to\mathcal X$ is finite and by proper base change, one see $X_0\to X'$ is an isomorphism onto a component of $X'$. So if $t\notin \mathcal Y\cup\mathcal W\cup\mathcal V$, $X_0$ is a component of $X'$ and hence geometrically normal. Otherwise $t\in\mathcal V':=\mathcal V\cup\mathcal W\cup\mathcal Y$, we replace $\mathcal T$ by $\mathcal V'$ and $\mathcal X$ by $\mathcal X\times_{\mathcal T}\mathcal V'$ and do the same arguments above, by Noetherian induction we are done.
\end{proof}

\paragraph{\tbf{Openness of klt locus}}
We show that klt locus for Fano type log Calabi-Yau pairs are open in sufficiently large linear systems.
\begin{lemma}[Openness of klt locus for Fano type threefolds]
\label{kltopen}
    Let $X$ be a ${\mathbb Q}$-factorial Fano type threefold in char $k>5$, let $(X,\Delta)$ be a klt pair with $K_X+\Delta\sim_{\mathbb Q}0$. Then for some large $n$ such that $n\Delta$ is integral and an open neighbourhood $V\subset |n\Delta|$ of $n\Delta$, for any $L\in V$, $(X,\frac{1}{n}L)$ is klt.
\end{lemma}
\begin{proof}
    We run $-K_X$-MMP on $X$, we end with a model $(X',\Delta')$. Here $K_{X'}+\Delta'\sim_{\mathbb Q}0$ and $f:(X,\Delta)\dashrightarrow(X',\Delta')$ is crepant. We will have $\Delta'=f_*\Delta$. Now for any $0\leq D'\sim_{\mathbb Q}\Delta'$ on $X'$, take a common resolution of $X$ and $X'$, say $\phi:W\to X$ and $\phi':W\to X'$. Set $D:=f^\sharp D':=\phi_*\phi'^*(K_{X'}+D')-K_X=:\phi_*(K_W+D_W)-K_X$. Since the MMP is $K_X$-positive, $D$ is effective by negativity lemma. We claim that
    \[f_*:|\Delta|_{\mathbb Q}\leftrightarrows|\Delta'|_{\mathbb Q}:f^\sharp\] is a bijection. First we prove $f_*f^\sharp=id$, consider the crepant maps \[(X,D)\leftarrow (W,D_W)\to(X',D').\]
    We have $f_*D=f_*(K_X+D)-f_*K_X=\phi'_*(K_W+D_W)-K_{X'}=D'$ as desired. Now we only need to show $f_*$ is injective. Indeed if $f_*D_1=f_* D_2=D'$, then $\phi^*(K_X+D_1)=\phi^*(K_X+D_2)=\phi'^*(K_{X'}+D')$, which means that $\phi^*D_1=\phi^*D_2$ and hence $D_1=D_2$. Moreover, we see that if $nD_1\sim nD_2$, then $n(K_X+D_1)\sim n(K_X+D_2)$ and $n(K_{X'}+D'_1)\sim n(K_{X'}+D'_2)$ and $n D'_1\sim nD'_2$. So $|n\Delta|\simeq |n\Delta'|$ as varieties. For any $D\in \frac{1}{n}|n\Delta|$, $(X,D)\dashrightarrow (X',D')$ is crepant, thus it suffices to prove the klt property for $(X,D')$. Since $\Delta'$ is nef and big, the assertion follows from Bertini's theorem for hyperplane sections in arbitrary characteristic and resolution of singularities \cite{BertiniHyperplane}.
\end{proof}



\section{Adjunctions}












%%%%%%%%%%%%%%%%%%%%%%%%%%%%%%%%
In this chapter, we introduce the adjunction formulas, which are crucial tools for inductions. In general, adjunction formulas relates the log canonical divisors of two varieties. We will introduce the divisorial adjunction and the canonical bundle formula here.

\paragraph{\tbf{Divisorial adjunction}}
The divisorial adjunction relates the singularity of a pair and its restriction on a component of the boundary. Explicitly speaking, we have:
\begin{theorem}[\cite{birkar2014existenceflipsminimalmodels} 4.1,4.2]\label{divadj}
    Let $(X,B)$ be a pair, $S$ be a component of $\lfloor B\rfloor$, and $S^\nu\to S$ be the normalization. Then there is a canonically determined $\mathbb R$-divisor $B_{S^\nu} \geq 0$ such that
\[K_{S^\nu} + B_{S_\nu} \sim_{\mathbb Q} (K_X + B)|_S\]
Moreover let $\Phi\subset [0, 1]$ be a DCC set of rational numbers, assume that:
\begin{enumerate}
    \item $(X,B)$ is lc outside a codimension 3 closed subset, and
    \item the coefficients of $B$ are in $\Phi$.
\end{enumerate}
Then $B_{S^\nu}$ is a boundary with coefficients in $\mathfrak S_\Phi$, here
\[\mathfrak S_\Phi=\{\frac{m-1}{m}+\sum\frac{l_ib_i}{m}\leq 1|m\in\mathbb Z^{>0}\cup\{\infty\},l_i\in\mathbb Z^{\geq 0},b_i\in\Phi\}.\]

\end{theorem}

Suppose $B\in\Phi(R)$ for some finite set $R\subset[0,1]\cap {\mathbb Q}$, let $I$ be a common denominator of $R$, say $R=\{\frac{r_i}{I}|0\leq r_i\leq I\}_i$ and choose any element $\alpha\in \mathfrak S_{\Phi(R)}$, say
\[\alpha=\frac{m-1}{m}+\sum_{i,n}\frac{l_{i,n}(1-\frac{r_i}{In})}{m}\leq 1.\]
Then we have $\sum\limits_{i,n}l_{i,n}(1-\frac{r_i}{In})\leq 1$. Let $r$ to be the biggest one in all $r_i$'s. If $l_{i,n}\neq 0$ for some $i$ and $n>1$, and if some $l_{i,n}\neq 0$ for another $i,n$ and $1-\frac{r_i}{In}\neq 0$, then we have \[1-\frac{r}{In}+1-\frac{r}{I}\leq 1,\] which implies
\[\frac{r}{I}(1+\frac{1}{n})\geq 1.\]
Hence if $n>I$, then we will get a contradiction. Let $n_\alpha$ to be the largest $n$ such that $l_{i,n}\neq 0$, then for $n_\alpha>I$, we have
\[\alpha=\frac{m-1}{m}+\frac{(1-\frac{r_i}{In})}{m}=1-\frac{r_i}{mIn}\in\Phi(R).\]
If $n_\alpha>1$, then there is only one $l_{i,n}\neq 0$ for $n>0$, which is exactly equal to $1$. we have that
\[\alpha=\frac{m-1}{m}+\frac{1-\frac{r_i}{In}}{m}+\sum_j\frac{l_j(1-\frac{r_j}{I})}{m}=1-\frac{r_i}{mIn}+\sum_j\frac{nl_j(1-\frac{r_j}{I})}{mn}\leq 1\]
which requires that \[\sum_j nl_j(I-r_j)\leq r_i.\]
hence there are only finitely many choices of the set $\{l_j\}|_{r_j\neq I}$ for finitely choices $r_i$ and finitely many choices $n$. So there are only finitely many choices for $\frac{r_j}{In}+\sum\limits_jl_j(1-\frac{r_j}{I})$. Add them to $R$ we get a new finite set $R'$, and $\alpha\in\Phi(R')$ by construction. Similar story happens for $n_\alpha=1$ and hence there is a finite set $S$ depending only on $R$ such that $\alpha\in\mathfrak S_{\Phi(R)}\subset \Phi(S)$.

Divisorial adjunction also closely relates the singularities of the pairs involved. We have the theorem of adjunction and inversion of adjunction in $\chara>5$.

\begin{theorem}[\cite{threedimmmp}, Theorem 6.2]\label{divadj2}
    Let $(X, S + B)$ be a pair where $X$ is a 3-dimensional normal variety in characteristic $p > 5$, $S+B\geq 0$ is a ${\mathbb Q}$-divisor such that $\lfloor S + B\rfloor = S$ is a prime Weil divisor. Let $K_{S^\nu}+B_{S^\nu}=(K_X+B+S)|_{S^\nu}$ to be the divisorial adjunction, then
    \begin{enumerate}
        \item $(X,S+B)$ is lc on a neighborhood of $S$ if and only if $(S^\nu,B_{S^\nu})$ is lc.
        \item If $X$ is ${\mathbb Q}$-factorial, then $(X,S+B)$ is plt on a neighborhood of $S$ if and only if $(S^\nu, B_{S^\nu})$ is klt. Moreover if $(X,S+B)$ is plt, then $S$ is normal.
    \end{enumerate}
\end{theorem}
If we consider F-singularities, we will have a similar theory of F-adjunctions, cf. \cite{Das_2016}. 

\paragraph{\tbf{Canonical bundle formula}}
Another useful adjunction is the fiber space adjunction, or more usually called the canonical bundle formula:
\begin{conjecture}[Canonical bundle formula]
\label{cbff}
    Suppose $(X,B)/Z$ is a contraction, where $(X,B)$ is a generically lc projective pair and $K_X+B\sim_{\mathbb Q} 0/Z$, let $\eta=\text{Spec}(k(D))$ be the generic point of any given prime divisor $D$ on $Z$, then there is 
    \[B_Z:=\sum\limits_{D\text{ prime divisor on }Z}(1-lct_\eta(X,B,f^*D))D\]
     and some pseudo-effective ${\mathbb Q}$-b-divisor $M_Z$ such that $K_X+B_X\sim_{\mathbb Q} f^*(K_Z+B_Z+M_Z)$. Moreover, if $(X,B)$ is lc, then $M_X$ is a b-nef ${\mathbb Q}$-b-divisor, where $M_X=K_X+B_X-f^*(K_Z+B_Z)\sim_{\mathbb Q}f^*M_Z$
\end{conjecture}
Here, $B_Z$ is called the determinant part, which is determined uniquely, and $M_X$ is called the moduli divisor. As $M_X=f^*M_Z$, we also call $M_Z$ the moduli part of the fibration. As $f:X\to Z$ is a contraction, the b-nefness of $M_X$ is equivalent to that of $M_Z$. 

In characteristic $0$, the conjecture was completely resolved by a series of works (cf.\cite{Kawamata1997SubadjunctionOL} \cite{Ambro_2005}\cite{2021arXivpositivityofmoduli}\cite{FilsemiampleModulipart}\cite{JLX22}\cite{CHLX23}.) In positive characteristic, however, the conjecture fails in general (cf. \cite[3.5]{Witaszek2017OnTC}).  Some works \cite{Witaszek2017OnTC}\cite{benozzocanonicalbundleformula} have resolved the case when the geometric generic fiber, or equivalently, the general fibers are lc and the base or the fibers are projective curves. In \cite{jxtcbf}, we prove that for Fano type fibration of a threefold with normal general fibers, the canonical bundle formula holds in large characteristics. Using \Cref{genorm} and the similar method in \cite{jxtcbf}, we can remove the normality condition. Explicitly speaking, we have:

\begin{theorem}[Canonical bundle formula for Fano type fibrations]
\label{cbf1}
    Assume $(X,B)$ is a lc pair over an algebraically closed field of characteristic $p$, $f:X\to Z$ is a contraction with dim $X=3$ and dim $Z>0$, $K_X+B\sim_{\mathbb Q} 0/Z$, $B$ is relatively big. Let $\epsilon>0$ be a real number and $R\subset [0,1]$ be a finite set of rational numbers, then there is a prime number $p_0=p_0(\epsilon,R)$ such that if $\chara k=p>p_0$ and one of the following conditions holds:
    \begin{enumerate}
        \item $\dim Z=3$, i.e. $X\to Z$ is birational,
        \item $\dim Z=2$, $B^{h}\in\Phi(R)$ where $B^h$ is the horizontal part of $B$,
        \item $\dim Z=2$, $(X,B)/Z$ has lc general fibers,
        \item $\dim Z=1$, $(X_\eta,B_\eta)$ is $\epsilon$-lc and $B^{h}\in\Phi(R)$,
        \item $\dim Z=1$, $X_\eta$ is an $\epsilon$-lc Fano variety and $B^h\in\Phi(R)$.
    \end{enumerate}
    Then we have the following formula:
    \[K_X+B\sim_{\mathbb Q} f^*(K_Z+B_Z+M_Z)\]
    where \[B_Z:=\sum_{D\text{ prime divisor on }Z}(1-\text{lct}_{\eta(D)}(f^*D,X,B))D\]
    is the discriminant part and $M_Z$ is a b-nef ${\mathbb Q}$-b-divisor.
\end{theorem}
\begin{proof}
    The proof is almost the same idea of \cite[3.3]{jxtcbf}, but we state it here systemically for convenience. The idea of the proof is to reduce to the case where the general fibers are lc and the base or the fiber are smooth curves, and then apply the known results.

    When $\dim Z=3$, in this case we have that \[K_X+B=f^*f_*(K_X+B)\sim_{\mathbb Q} f^*(K_Z+B_Z+M_Z).\] For a divisor $D$ on $Z$, we denote its birational transformation also by $D$ for convenience. Then
    \[
    \begin{aligned}
        \mu_D(B_Z)&=1-\lct_\eta(f^*D,X,B)=1-\sup\{t|(X,B+tD)\text{ is lc near $D$}\}\\&=1-(1-\mu_D(B))=\mu_D(B)=\mu_D(f_*(K_X+B)-K_Z).
    \end{aligned}\]
    Hence $M_Z=0$ and the result follows.
    
    When $\dim Z=2$. In this case, the general fibers are curves. We first reduce the condition 2 case to condition 3 case. That is, we are going to prove the general fibers are lc. Replacing $X$ by its resolution $X'$, we have a crepant model $(X',B'=B^+-B^-)\to (X,B)$, where $(X',B')$ is sub-lc. Also, replace $Z$ by its smooth locus, one may assume that $f:(X,B)\to Z$ is a fibration between smooth quasi-projective varieties. Moreover, the image of the exceptional divisors of $X'\to X$ under $f$ is not surjective on $Z$ since codim$(Sing(X)/X)\geq 2$, hence the horizontal parts of $B'$ and $B$ are the same.
        
    Consider $X_\eta$ the generic fiber of $X/Z$, then it is projective normal, Fano type and Gorenstein as $X$ is smooth. Suppose $\pi:Y\to X_\eta$ is a normalization of an irreducible component of $X_{\bar\eta}$, then by the behavior of canonical divisor under inseparable base changes \cite[1.1]{patakfalvi2018singularitiesgeneralfiberslmmp}, we have \[K_Y+(p-1)C\sim\pi^*K_{X_\eta}\] where $C$ is an integral effective divisor (not 0 iff $X_\eta$ is not geometrically normal), hence $K_Y\sim-(p-1)C+\pi^*K_{X_\eta}$ is anti-ample since deg$(-K_{X_\eta})>0$. Moreover \[-2\leq\text{deg}(K_Y)\leq\text{deg}((1-p)C)\leq 1-p,\] which implies $p\leq 3$, so we choose $p_0>3$.

    Geometric normality, geometric regularity and geometric reducedness is an open condition (cf. \cite[II.6.9.1 and III.12.1.1]{PMIHES_1967__32__5_0}). Since $X/Z$ admits normal geometric generic fiber and $k$ is algebraically closed, we have that a general fiber $F$ is a projective smooth curve. Consider $(F,B_F)$, if a component $D$ of $B$ is vertical, then its contribution in $B_F=B|_F$ is $0$ as $F$ is a general fiber. So $B|_F=B^{hor}|_F$ and we only need to consider the horizontal part which impacts the singularity of the general fiber. For a horizontal component $D\in Supp(B)$, we can move the general fiber such that $D|_F$ has no multiplicity except the inseparable part.
        
    Consider the surjection $D^\nu\to Z$ with inseparable degree $p^{k_D}$ of the field extension $K(D^\nu)/K(Z)$, then by the basic intersection theory, $D|_F=\sum p^{k_D} D_i$ where $D_i$ are reduced divisors. Hence, set $B^{hor}=\sum a_jD_j$ where $a_j\in \Phi$, we have \[0\sim_{\mathbb Q} (K_X+B)|_F=K_F+B^{hor}_F=K_F+\sum_j\sum_ip^{k_{D_j}}a_jD_{j,i}.\] If some $k_{D_j}>0$, then calculating the degree on both sides, one sees that \[0\geq \text{deg}(K_F)+p\min(\Phi)\geq-2+p\min(\Phi),\] hence $p\leq \frac{2}{\min(\Phi)}$. Set $p_0\geq\frac{2}{\min(\Phi)}$ there is no inseparable horizontal part and the general fibers are lc. Hence we fall in the condition 3 case. By \cite{Witaszek2017OnTC} the canonical bundle formula holds in this case. 
        
    In fact,  we have proved that the geometric generic fibre $X_{\bar{\eta}}$ is a smooth rational curve since it is normal and has a canonical divisor with a negative degree. Moreover we have proved that the horizontal part of $B$ is separable over $Z$ under condition 2, hence they are given by a union of rational sections of $X/Z$. We can use the arguments in the proof of \cite[3.1]{Witaszek2017OnTC} to see that $M_Z$ is semi-ample outside a codimension-$2$ closed subset, hence semi-ample since $Z$ is a surface. For condition 3 case, remove the isolated points of $Z$ where the horizontal part of $B$ does not extend and the singular locus of $Z$ away from $Z$, consider the morphism $Z\to \overline{M_{0,m}}$ defined by map $z$ to the $(X_z,\Supp(B_z))$. Then we see that $(X_z,\Supp(B_z))$ is lc for general $z$. Thus the inversion of adjunction for $(U,B^h_U+\phi_U^{-1}Q)$ in the proof of \cite[3.1]{Witaszek2017OnTC} would also apply and the rest proof would follow similarly as under the condition 2 case.
        
    When dim $Z=1$, the general fibers are now surfaces, and from now on one assumes $p>5$. Take a ${\mathbb Q}$-factorialization, we may assume $(X,B)$ is dlt ${\mathbb Q}$-factorial. As $X\to Z$ is a contraction, we have that the generic fiber $X_\eta$ is defined on $\eta=k(Z)$ and $H^0(X,\mathcal O_X)=\eta$. We reduce the condition 4 case to the condition 5 case. Firstly, after running $-(K_X+(1-\upsilon)B)\sim_{\mathbb Q}\upsilon B$-MMP on $(X,B)/Z$ for some $\upsilon$ close to $0$, we end with a good minimal model $(X',B')$ by LMMP and $(X',B')$ is generically lc. Suppose $X'\to Y/Z$ is the semiample fibration with respect to $B'$, then one sees that \[h:(X,B)\to (Y,B_Y)/Z\] is crepant and $-K_Y\sim_{\mathbb Q}B_Y/Z$ is ample$/Z$ and $B_Y=h_*B$. Since $h$ is crepant$/Z$, the construction of the canonical bundle formula is compatible and one can assume that $B$ is ample$/Z$.

    After the reduction, the generic fiber is an $\epsilon$-lc del Pezzo surface and hence belongs to a bounded family over $\mathbb Z$ by \cite{bernasconi2024boundinggeometricallyintegraldel}. By \Cref{genorm}, we have that for $p>p_0$ large enough, the geometric generic fiber of $X/Z$ is normal, so is the general fibers. Since the coefficients of $B^h$ falls in a hyperstandard  set which is DCC, and $K_{X_\eta}+B_\eta\sim_{\mathbb Q} 0$, we have that $B^h$ falls in a finite set by \Cref{singgen}. As a result  $(X_\eta,B_\eta)$ falls in a log bounded set$/\Spec \mathbb Z$ by Hilbert scheme arguments.

    Now there is a very ample divisor $H$ on $X_\eta$ with $-K_X\cdot H=B\cdot H$ bounded and $H^2$ bounded by some natural number $M$. Moreover $H$ induces a closed immersion $X_\eta\to \mathbb P^N_\eta$, hence its base change to algebraic closure $\bar H$ will also induce a closed immersion $X_{\bar\eta}\to \mathbb P^N_{\bar\eta}$, which implies that the geometric generic fiber also falls in a bounded set over $Spec\mathbb Z$. Consider the $\eta$-non-smooth (closed since $X_\eta$ is geometrically normal) points on $X_\eta$ and the non-snc points of $B_\eta$, we call them $x_i$'s. As $X_\eta$ is a klt surface, we see $X_\eta$ is ${\mathbb Q}$-factorial. Moreover, the Cartier index of $K_{X_\eta}$ and $B_\eta$ near $x_i$ is bounded by some natural number $N$ and the defining equations of $X_\eta$ have bounded degrees $N$ since $(X_\eta,B_\eta)/k$ falls in a bounded set$/Spec\mathbb Z$. Thus, for the closed points $x_i$, $\kappa(x_i)/\eta$ has a degree bounded by $N$. Hence, for $p>N$, $\kappa(x_i)/\eta$ is separable. Now we are going to show that the pair $(X_\eta,B_\eta)$ is geometrically lc if $B_\eta$ is geometrically reduced.

    In the following two paragraphs, we denote $k$ to be $\eta$ and denote $X$ to be $X_\eta$ for convenience and we are already over an imperfect base field, the following arguments are similar to those in \cite{sato2024generalhyperplanesectionslog}. In fact, after shrinking $X$, we only need to consider a surface singularity $(x\in X,B)$, here $x=\Spec k\to X$ is a closed point. If $x\in X$ is smooth and $B$ is snc near $x$, since $\kappa(E)/k$ is smooth for each component $E$ of $B$, then the components of $B$ and the intersections are smooth over $\kappa(x)$, and therefore $(x\in X,B)$ is geometrically log smooth and the discrepancies will remain the same after base change to geometric case. So we may assume that $(x\in X,B)$ is an isolated (geometric) surface singularity. Pick a log resolution of singularity $f:Y\to X$ near $x$, say Exc$(x)=\sum\limits_iE_i$ is the sum of the exceptional divisors on $Y$, which is a scheme defined over $\kappa(x)$. Suppose $l$ is a purely inseparable field extension of $\kappa(x)$, then for any irreducible $\kappa(x)$-scheme $X$, $X\times_{\kappa(x)}l$ is homeomorphic to $X$ as topological spaces and hence also irreducible. Hence, we can take a finite separable field extension $k'$ of $\kappa(x)$ such that the irreducible components of $E_i\times_{\kappa(x)}k'$ are geometrically irreducible. Since $k'/\kappa$ is finite separable, and hence a single extension by a monic polynomial $g(t)\in \kappa[t]$. Suppose $X=Spec(R)$ with $\kappa=R/\mathfrak m$, take a monic lift $\tilde g(t)\in R[t]$ of $g(t)$, we have that \[X'=Spec(R[t]/\tilde g(t))\to Spec(R)=X\] is finite étale and surjective, with the corresponding $x'\to x$ realizing $\kappa\to k'$.  Then $Y'=Y\times_XX'$ is a minimal log resolution of singularity by étale descent, whose components of exceptional locus are exactly such components of $E_i\times_\kappa k'$ and hence geometrically irreducible. Also $X'$ is klt with the same discrepancies on certain components by étale descent, and hence has rational singularities. Hence, we have $g(E_i)=0$ and by the classification of the dual graph \cite[A.3]{sato2024generalhyperplanesectionslog}, we have $\dim_{\kappa}H^0(E_i,\mathcal O_{E_i})\leq 4$. Hence $K=H^0(E,\mathcal O_E)$ is separable over $\kappa(x)$ and hence over $k$ if $p>3$, which means that $E$ is geometrically reduced as $E$ is reduced$/K$ which can be realized as a conic in $\mathbb P_{K}^2$ and no purely inseparable field extension $L/K$ will make $E_L$ to be a $p$-power multiple of the certain divisor since $E_L=\mathcal O_{\mathbb P_L^2}(2)$ and $p>2$.

    Now $E_i$ are geometrically integral and their base change to $\bar \kappa(x)$ are integral curves with arithmetic genus $0$, hence $\mathbb P^1$, which are smooth. Moreover, we have $B_i$'s that are geometrically reduced (and hence integral by a finite étale base change) by the assumption, hence the exceptional locus together with the strict transform of $B$ is snc and consists of geometrically integral regular divisors, moreover the exceptional divisors are smooth. Consider each $B_i$, which is a geometrically integral regular curve over $k$. If some $B_i$ is not geometrically normal, denote $C_i$ to be the normalization of its geometric model, then we have for some integral divisor $D$ on $C_i$ and\[K_{C_i}+(p-1)D\sim_{\mathbb Q} K_{B_i}|_{C_i}.\] Count the degree, since $B_i$ is Gorenstein and by Riemann-Roch, we have \[2g_C-2+(p-1)deg(D)=2g_{B_i}-2\] and hence $g_{B_i}\geq\frac{p-1}{2}$. However, in our case $B_i$ is a component of the boundary of some log bounded family $(X,B)$, which means $g_{B_i}$ is bounded by some $G$, hence this cannot happen if $p>2G+2$ when $D\neq 0$. Hence,  assume $p>2G+2$, all $E_i$'s and $B_i$'s are smooth$/k$. View $Exc(x)+\tilde B$ as a scheme over $x=\Spec k$, then all the extension degrees of $H^0(F,\mathcal O_F)/k$ of any component $F\subset Exc(x)+\tilde B$ are bounded since $(X,B)$ belongs to a log bounded family$/\mathbb Z$, which means the defining equation of $(x\in X,B)$ has bounded degree, hence when $p$ is large enough, all intersections $E_i\cap E_j$ or $E_i\cap B_j$ are separable over $x\in X$. Hence $(Y,Exc(x)+\tilde B)$  is snc after arbitrary field base change and we see that all the log discrepancies would remain the same after base change to the algebraic closure and the geometrically.
        
    Returning to our initial fibration setting, we now show that $B_{\eta}$ is geometrically reduced. Suppose $B=B^{hor}+B^{ver}=\sum a_jD_j+B^{ver}$, then the vertical part does not affect the singularities of the general fibers. After restricting to the geometric generic fiber $F$ we have \[0\sim_{\mathbb Q} (K_X+B)|_F=K_F+\sum_ja_{j}p^{k_{D_j}}D_{j,F}^{red}\] where $p^{k_{D_j}}$ is the inseparable degree of (the Stein factorization of) the map $D_j\to Z$. Consider  the bounded flat family $(\mathcal X,\mathcal B)/\mathcal T$ over $\Spec \mathbb Z$ which contains $(X_\eta,B_\eta)$, we see $(X_\eta,B_\eta)$ is a fiber contained in $(\mathcal X_\eta,\mathcal B_\eta)/\mathcal T_\eta$, which is the fiber of $(\mathcal X,\mathcal B)/\mathcal T$ over the unique map $\eta\to\Spec \mathbb Z$. Hence, $(F,B_F)$ is a fiber of $(\mathcal X_{\bar\eta},\mathcal B_{\bar\eta})/\mathcal T_{\bar\eta}$, which also falls in the bounded family$/\Spec \mathbb Z$. $F$ is normal and Fano by our assumption. Apply our arguments above to $(X_\eta, 0)$, we see $F$ is also $\epsilon$-lc and hence $\mathbb Q$-factorial. Then since $(F,B_F)$ falls into a log bounded set, we see the isolated singularities $x\in F$ are of bounded type$/\Spec\mathbb Z$ by considering the local defining equations when $p$ is large enough. Since $F$ is $\mathbb Q$-factorial, the local Cartier index at $x\in F$ is bounded by some sufficiently large integer $N$. As a result, we have $\bar H\cdot D_{j,F}^\text{red}\geq \frac{1}{N!}$ as $N!D_{j,F}^\text{red}$ is Cartier, here $\bar H$ is the very ample. Also, we have $-\bar H\cdot K_F\leq M$ for some sufficiently large integer $M$. If there is some $k_{D_j}>0$, then \[M\geq -K_F\cdot \bar H\geq pa_jD_{j,F}^{red}\cdot\bar H\geq \frac{p\min(\Phi)}{N!},\]leading to a contradiction when $p>\frac{N!M}{\min(\Phi)}$. Thus when $p$ is large enough, we see there is no inseparable part both in the map from the exceptional divisors and the components of $B$ to $Z$, in particular the log discrpancies are kept and hence $(F,B_F)$ is geometrically lc and by \cite[0.2]{benozzocanonicalbundleformula}, the canonical bundle formula holds in this case.
\end{proof}

A direct corollary of the canonical bundle formula is that a contraction of a Fano type threefold is of Fano type if the base is not an elliptic curve (though the author believes  that this will not happen when the characteristic is large enough under the conditions of the following corollary). Explicitly speaking, we have:
\begin{corollary}
\label{contr}
       Assume $(X,B)$ be a $3$-dimensional projective lc pair with a contraction $f:X\to Z$, where $K_X+B\sim_{{\mathbb Q}}0/Z$. Let $R\subset [0,1]$ be a finite set of rational numbers and $\epsilon>0$ be a real number, then there exists a prime number $p_0$ depending only on $R$ and $\epsilon$ such that if $-(K_X+B)$ is nef, $X$ is of Fano type, char $k=p>p_0$, $B\in\Phi(R)$ and one of the following conditions holds:
    \begin{enumerate}
        \item $Z$ is a projective normal surface,
        \item $Z$ is a projective normal curve, $(X,B)$ is $\mathbb Q$-factorial dlt and $X_\eta$ is an $\epsilon$-lc Fano surface.
    \end{enumerate}
    Then $Z$ is of Fano type, except the only case when $X\to Z$ is a surface fibration over an elliptic curve and the canonical bundle formula is written as
    \[K_X+B\sim_\mathbb Qf^*K_Z\sim_\mathbb Q0.\]
    Here $B_Z=M_Z=0$, $(X,B)$ is not klt and not geometrically klt, $B^\text{ver}=0$ and each normalization $S^\nu$ of a lc place $S\in\lfloor B\rfloor$ is not of Fano type.
\end{corollary}
\begin{proof}
    Taking a dlt ${\mathbb Q}$-factorialization of $X$, we may assume $X$ is dlt ${\mathbb Q}$-factorial. When $Z$ is a surface, we have that the horizontal part of $B$ are separable over $Z$ when $p>p_0$ and $(X,B)$ has lc general fibers. $-(K_X+B)$ is semi-ample since $X$ is of Fano type. Hence if $-(K_X+B)\sim_{\mathbb Q}f^*D$ for some $D$ on $Z$, then $D$ is semi-ample. Let $(X,\Delta)$ be a klt pair such that $K_X+\Delta\sim_{\mathbb Q}0$, we see that $\Delta$ is big. Since the klt condition is an open condition \Cref{kltopen}, that is, for some sufficiently large $n$ such that $n\Delta$ is integral, there exists an open neighbourhood $V\subset |n\Delta|$ of $n\Delta$ such that for any $L\in V$, the pair $(X,\frac{1}{n}L)$ is klt.  Since $V$ is open, the base locus of $V$ is equal to the base locus of $|n\Delta|$. Suppose $E$ is a divisorial part of $\text{Bs}_\mathbb Q(\Delta)$, then $E$ is also a divisorial part of $\text{Bs}_\mathbb Q(B)$ since $\Delta-B=(K_X+\Delta)-(K_X+B)\sim_\mathbb Q-(K_X+B)$ is semi-ample. Hence we can enlarge $n$ to be sufficiently divisible such that $n\Delta\sim n(F+M)$ for some $F$ fixed and $M$ movable in $V$ and moreover $F_\eta\leq B_\eta$. In particular, when $n\Delta'\sim n\Delta$ varies inside $V$, $M_\eta$ varies base-point freely on the generic fiber. Hence there is some $M'$ induced by some $n\Delta'\in V$ such that \[\begin{aligned}(X_\eta,(1-\epsilon)B_\eta+\epsilon \Delta'_\eta)&=(X_\eta, (1-\epsilon)F_\eta+\epsilon F_\eta+(1-\epsilon)(B_\eta-F_\eta)+\epsilon M'_\eta)\\&=(X_\eta,B_\eta-\epsilon(B_\eta-F_\eta)+\epsilon M'_\eta)\end{aligned}\] is geometrically lc for sufficiently small $\epsilon$. Moreover we see that if the lc place $\lfloor B_\eta\rfloor\leq F_\eta$, then $\lfloor F\rfloor\neq 0$ which is a contradiction as $F\leq \Delta$ which has zero integral part. Hence $(X_\eta,B_\eta-\epsilon(B_\eta-F_\eta)+\epsilon M'_\eta)$ is klt for sufficiently small $\epsilon$. As $B^{\text{hor}}$ is separable over $Z$, so if $D\in \lfloor B_{\bar\eta}\rfloor$, then there exists some $D'\in \lfloor B_{\eta}\rfloor$ whose base change to the algebraic closure contains $D$. Consequently, $(X_\eta,B_\eta-\epsilon(B_\eta-F_\eta)+\epsilon M'_\eta)$ is geometrically klt.

    Now consider $\Gamma=(1-\epsilon)B+\epsilon\Delta'$ for some $\epsilon\ll 1$, we have that $(X,\Gamma)$ is klt and geometrically generically klt. Let $A$ be an ample divisor on $Z$, we see that for some $\zeta\ll 1$, there is an ample ${\mathbb Q}$-divisor $H\geq \zeta f^*A$ and some $\Delta''\sim_{\mathbb Q} \Delta$ such that $\Delta''=D+H$ for some effective ${\mathbb Q}$-divisor $D$. Thus take $\delta\ll 1$, set \[\Gamma':=(1-\epsilon)B+\epsilon(1-\delta)\Delta'+\epsilon\delta D+\epsilon\delta(H-\zeta f^*A)\sim_{\mathbb Q}\Gamma-\delta\epsilon\zeta f^*A.\]
    We have that $(X,\Gamma')$ is klt and geometrically generically klt when $\delta$ is small enough. Moreover, \[K_X+\Gamma'\sim_{{\mathbb Q},Z}(1-\epsilon)(K_X+B)+\epsilon(K_X+\Delta)\sim_{{\mathbb Q},Z}0.\]
    Hence the canonical bundle formula for curve fiber case would apply, that is we have some $b$-nef ${\mathbb Q}$-$b$-divisor $M_Z$ such that:
    \[K_X+\Gamma'\sim_{\mathbb Q}f^*(K_Z+\Gamma'_Z+M_Z)\sim_{\mathbb Q} -(1-\epsilon)f^*L-\delta\epsilon\zeta f^*A\]
    Thus take some $\eta\ll 1$, we have that 
    \[-(K_Z+\Gamma'_Z+M_Z+(1-\eta)\delta\epsilon\zeta A)\sim_{\mathbb Q}(1-\epsilon)L+\delta\epsilon\zeta\eta A.\]
    By construction, $(Z,\Gamma'_Z)$ is klt and $M_Z+(1-\eta)\delta\epsilon\zeta A$ is nef and big on some higher model. Thus passing to high resolutions of $X/Z$, there is some $\Omega\sim_{\mathbb Q}\Gamma'_Z+M_Z+(1-\eta)A$ on $Z$ such that $(Z,\Omega)$ is klt and $-(K_Z+\Omega)$ is nef and big since $L$ is semi-ample and $A$ is ample. In particular, $Z$ is of Fano type.

    When $Z$ is a curve, we first prove the case when $(X,B)$ is $\mathbb Q$-factorial dlt and $(X,B)/Z$ is of relative Picard number $1$. by the proof of \Cref{cbf1}, we see $(X,B)$ is geometrically generically lc and $B_\eta$ is geometrically reduced. If $(X,B)$ is geometrically generically klt, then the same construction in curve fiber case for $\Gamma=(1-\epsilon)B+\epsilon \Delta'$ and $\Gamma'$ would applies as $(X,\Gamma)$ is klt and geometrically generically klt for any $\Delta'\in V$ and arbitrary sufficiently small $\epsilon$. Moreover, by \cite[0.2]{benozzocanonicalbundleformula}, the canonical bundle formula for $(X,\Gamma')/Z$ will also hold and the same proof as the curve fiber case would apply. So we only need to prove the case when $(X,B)$ is not geometrically klt. Since $(X,B)$ is $\mathbb Q$-factorial dlt, there is some component $S\subset\lfloor B\rfloor$ which is horizontal$/Z$ and $S_\eta$ is a lc center in $(X_\eta,B_\eta)$. Since the finite part of the stein factorization of $S\to Z$ contains no inseparable part, by connectedness principle \cite[1.2]{YH17wittnadelvanishing}, the morphism $S\to Z$ is a contraction. Consider $(S^\nu,B_{S^\nu})/Z$ given by divisorial adjunction,  then $-(K_{S^\nu}+B_{S^\nu})=-(K_X+B)|_{S^\nu}$ is nef and relatively trivial over $Z$, and $B_{S^\nu}\in\Phi(J)$ for some finite set depending only on $R$. If $S^\nu$ is of Fano type, then consider the pair $(S^\nu,B_{S^\nu})/Z$ and apply the same method in the curve fiber case and we are done. So we assume $S$ is not of Fano type from now on.  
    
    Suppose $Z$ is not a curve of Fano type, i.e $Z\neq \mathbb P^1$ for contradiction. Then by the canonical bundle formula applied to $(X,B)/Z$, we see that $K_X+B\sim_\mathbb Qf^*(K_Z+B_Z+M_Z)$ which is $\mathbb Q$-linearly effective, however $-(K_X+B)$ is semi-ample, we must have $B_Z=M_Z=0$, $Z$ is an elliptic curve and $K_X+B\sim_\mathbb Q0$. $B_Z=0$ implies that $f^*z$ has multiplicity $1$ for every $z\in Z$ and that $B^\text{ver}=0$.
\end{proof}

\section{Boundedness of complements for a lc-trivial fibration}











%%%%%%%%%%%%%%%%%%%%%%%%%%%%%%%%%%%%%%%%%%%%%%%%%%%%%%%%%%%%%%%%%%%%%%%%%%%%%%%%%%%%%%%
In this chapter, we prove that the relative complements of Fano type fibrations of dimension $3$ admits bounded complements over any point on the base. We first prove that the boundedness of complements can be deduced from sub-crepant birational models.
    \begin{lemma}
    \label{subcrep}
        Let $(X',B' + M')$ with data $\phi:X\to X'$ and $M$ on $X$, and $(X',B'' + M'')$ be 2 generalized pairs. Assume (by replacing $X$ with a higher resolution), $\phi:X\to X',\psi:X\to X''$ be a common resolution such that $\psi_*M = M''$. Suppose further that \[\phi^*(K_{X'} + B' +M') + P = \psi^*(K_{X''} + B'' +M'')\] for some $P \geq 0$, then if $(X'',B'' +M'')$ has an $n$-complement, then so does $(X',B' +M')$.
    \end{lemma}
    \begin{proof}
        Let $B''^+\geq B''$ be an $n$-complement for $(X'',B'' + M'')$. Consider \[B'^+ := B' + \phi_*(P + \psi^*(B''^+-B'')),\] then we get \[\begin{aligned}&n\phi^*(K_{X'} + B'^+ +M') = n\phi^*(K_{X'} + B' +M'+\phi_*(P + \psi^*(B''^+- B'') )\\&=n\phi^*(K_{X'} + B' +M'+\phi_*(\psi^*(K_{X''} + B'' +M'') + \psi^*(B''^+- B'')- \phi^*(K_{X'} + B' +M'))\\&=n(\phi^*(K_{X'} + B' +M')-\phi^*\phi_*\phi^*(K_{X'} + B' +M')+\phi^*\phi_*\psi^*(K_{X''} + B''^+ +M''))\\&=n\phi^*\phi_*\psi^*(K_{X''} + B''^+ +M'')\sim 0\end{aligned}\] by assumption. Hence, we obtain $n(K_{X'} + B'^+ +M')\sim0$. Moreover, $(X, B'^+ +M')$ is generalized lc  beacuase $(X'',B''^+ +M'')$ is generalized lc and $X$ is a common log resolution.
    \end{proof}
This lemma implies that the boundedness of complements is preserved under some important operations. If $(Y,B_Y+M_Y)\to (X,B+M)$ is a dlt ${\mathbb Q}$-factorialization, and if  $K_Y+B_Y+M_Y$ admits an $n$-complement, then $K_X+B+M$ will also admit an $n$-complement. If $X\dashrightarrow  X'$ is a partial step of $-(K_X+B+M)$-MMP, and if $K_{X'}+B'+M'$ admits an $n$-complement, then  $K_X+B+M$ will also admit an $n$-complement.

Now we give the lemma of effective canonical bundle formulas for Fano type fibrations which admit a bounded relative complement.
\begin{lemma}
\label{effcbf}
    Let $ R\subset [0,1]$ be a finite set of rational numbers, then there is $p_0>5,q\in\mathbb N$ and a finite set of rationals $ S\subset[0,1]$ depending only on $ R$ and satisfying the following. Assume $(X,B)$ is a pair and $f:X\to Z$ is a contraction such that:
    \begin{enumerate}
        \item $(X,B)$ is a projective lc pair of dimension $3$ over an algebraically closed field with char $p>p_0$,
        \item  $B\in\Phi( R)$,
        \item $K_X+B\sim_{\mathbb Q} 0/Z$,
        \item $X$ is of Fano type over $Z$,
        \item the canonical bundle formula \Cref{cbff} holds for $(X,B)/Z$,
        \item \Cref{cplt} holds for $(X,B)/Z$.
    \end{enumerate}
    Then we have $B_Z\in \Phi( S)$ and $qM_Z$ is b-nef b-Cartier and the effective canonical bundle formula:
    \[q(K_X+B)\sim qf^*(K_Z+B_Z+M_Z).\]
\end{lemma}
\begin{proof}
    Let $q=n$ be the number given by \Cref{cplt} which depends only on $ R$. There a $q$-complement $K_X+B^+$ of $K_X+B$ over some point $z\in Z$ with $B^+\geq B$. Since over $z$ we have $K_X+B\sim_{\mathbb Q}0$ and $q(K_X+B^+)\sim 0$, we have $B^+=B$ over $z$, therefore $q(K_X+B)\sim 0$ over the generic point of $Z$. Hence there is a rational function on $X$ such that $qL:=q(K_X+B)+div(\alpha)$ is zero over $\eta_Z$. Hence $L$ is vertical over $Z$ and $L\sim_{\mathbb Q} 0/Z$, we have $L=f^*L_Z$ for some $L_Z$ on $Z$, let the moduli part $M_Z:=L_Z-(K_Z+B_Z)$ where $B_Z$ is the discriminant part for $(X,B)/Z$. We have \[q(K_X+B)\sim qL=qf^*L_Z=qf^*(K_Z+B_Z+M_Z).\]

    In the following two paragraphs, we first reduce the proof that $B_{Z}\in\Phi(S)$ and that $qM_{Z}$ is integral to the case where the base is a curve. Assume dim $Z>1$, let $H$ be a general very ample divisor of $Z$ and $G$ its pullback to $X$, let \[K_G+B_G=(K_X+B+G)|_G.\] By adjunction $B_G\in\Phi(R')$ for some finite set $R'$ of rational numbers and $(G,B_G)$ is lc. Since $G\to H$ is at most with curve fibers and $G$ is of Fano type over $H$ by adjunction, the canonical bundle formula applies for $(G,B_G)/H$ by possibly enlarging $p_0$ by \Cref{cbf}. Let $g:G\to H$ be the induced map. Let $D$ be a prime divisor on $Z$ and $C$ be a component of $D\cap H$. Let $t:=\lct_{\eta(D)}(f^*D,X,B)$. Then there is a non-klt center of $(X,B+tf^*D)$ mapping onto $D$, which is also a non-klt center of $(X,B+G+tf^*D)$. Intersecting it with $G$ gives a non-klt center of $(G,B_G+tg^*C)$ mapping onto $C$ by inversion of adjunction. Thus $t=\lct_{\eta(C)}(g^*C,G,B_G)$, therefore $\mu_D(B_Z)=\mu_C(B_H)$. Since $B_G\in \Phi(R')$, by induction hypothesis, $B_H\in\Phi(S)$ for some finite set $S$ depending only on $R'$ and hence $R$. Therefore $B_Z\in\Phi(S)$.
    
    Now we are going to show $qM_Z$ is integral. Pick a general $H'\sim H$ and set $K_H:=K_Z+H'|_H$ as a Weil divisor, let $M_H:=(L_Z+H')|_H-(K_H+B_H)$ we have \[\begin{aligned}q(K_G+B_G)&\sim q(K_X+B+G)|_G\sim q(L+G)|_G\sim qf^*(K_Z+B_Z+M_Z+H')|_G\\&\sim qg^*(L_Z+H')|_H=qg^*(K_H+B_H+M_H)\end{aligned}\]
    Hence $M_H$ is the moduli part of $(G,B_G)/H$ and $B_H+M_H=(B_Z+M_Z)|_H$. Hence $\mu_C(B_H+M_H)=\mu_D(B_Z+M_Z)$, which means that $\mu_CM_H=\mu_DM_Z$ as $\mu_CB_H=\mu_DB_Z$. By induction hypothesis, $qM_H$ is integral, which means that $qM_Z$ is integral as desired. So we assume $Z$ is a projective normal curve now.
    
    We now show that $B_Z\in\Phi(S)$. Pick a closed point $z\in Z$, let $t=lct(X,B,f^*z)$ and $\Gamma=B+tf^*z$, let $(X',\Gamma')$ be a ${\mathbb Q}$-factorial dlt model of $(X,\Gamma)$ such that $\lfloor \Gamma'\rfloor$ has a component mapping to $z$. Then there is a boundary $B'\leq \Gamma'$ such that $B'\in\Phi( R)$ and $\lfloor B'\rfloor$ has a component map to $z$ and $B^\sim\leq B'$. Run MMP on $-(K_{X'}+B')\sim_{\mathbb Q} \Gamma'-B'$ over $Z$ and let $X''$ be the final model, then $X''$ is of Fano type over $Z$, $B^\sim\leq B''\in\Phi(R)$ and $-(K_{X''}+B'')$ is nef over $Z$ and $(X'',B'')$ is lc. Thus $K_{X''}+B''$ has a $q$-complement $K_{X''}+B''^+$ over $z$ with $B''^+\geq B''$. So $K_{X'}+B'$ admits a $q$-complement $(X',B'^+)$ with $B'^+\geq B'$ over $z$. Pushing $K_{X'}+B'^+$ down to $X$ we get a $q$-complement $K_X+B^+$ of $K_X+B$ over $z$ with $B^+\geq B$ such that $(X,B^+)$ has a non-klt center mapping to $z$. Since $K_X+B\sim_{\mathbb Q} 0/Z$ we have $B^+-B\sim_{\mathbb Q}0$ over $z$, hence $B^+-B$ is vertical effective over $Z$. Thus over $z$, $B'-B=sf^*z$. $s=t$ since a non-klt center of $(X,B^+)$ maps to $z$. We see the coefficient of $z$ in $B_Z$ is $1-t$, so we only need to show that $t$ falls in an ACC set. Let $S$ be a component of $f^*z$ and $b,b^+$ be its coefficients in $B$ and $B^+$. If $m$ is its multiplicity in $f^*z$, we have $b^+=b+tm$, so $t=\frac{b^+-b}{m}$. We see $b=1-\frac{r}{l}$ for some $r\in R$ and $l\in\mathbb N$, so $t=\frac{s}{m}$ where $s=b^+-1+\frac{r}{l}$. We see $b^+\leq 1$ and $qb^+$ is integral, we get \[1-\frac{1}{l}\leq b=1-\frac{r}{l}\leq b^+\leq 1-\frac{1}{q}\] or $b^+=1$, so $l\leq q$ and hence $s$ falls in a finite set of rational numbers, so $B_Z\in\Phi( S)$ for some finite set $ S$ of rational numbers.

    We now show that $qM_Z$ is integral. By $q(K_X+B)\sim qf^*L_Z$, we have $q(K_X+B)\sim 0$ over some non-empty open set $U\subset Z$ and $\text{Supp}(B_Z)\subset Z-U$. Set \[\Theta:=B+\sum\limits_{z\in Z-U}\text{lct}(X,B,f^*z)f^*z\] and $\Theta_Z$ to be the discriminant part of $(X,\Theta)/Z$, we have \[\Theta_Z=B_Z+\sum\limits_{z\in Z-U}\text{lct}(X,B,f^*z)z\] and Hence $\Theta_Z$ is a reduced divisor. Moreover we see that $K_X+\Theta$ is a $q$-complement of $K_X+B$ over each point $z\in Z-U$ by last step and hence $q(K_X+\Theta)\sim0/Z$. Therefore we have \[\begin{aligned}&q(K_X+\Theta)=q(K_X+B)+q(\Theta-B)\\&\sim qf^*(K_Z+B_Z+M_Z)+qf^*(\Theta_Z-B_Z)= qf^*(K_Z+\Theta_Z+M_Z)\end{aligned}.\] Since $q(K_{Z}+\Theta_{Z}+M_{Z})$ is Cartier and $K_{Z}+\Theta_{Z}$ is integral, it follows that $qM_{Z}$ is integral.

    In the following paragraphs, we show that $qM_Z$  is  a b-nef b-Cartier divisor. The b-nefness is ensured by the canonical bundle formula, so we only need to prove that $qM_Z$ is b-Cartier. We go back to general case that $Z$ may not be a curve and try to construct a birational model of $X'$ to make $qM_Z$ Cartier. Let $\phi:X'\to X$ be a log resolution of $(X,B)$ so that $X'\to Z'$ is a morphism and $Z'\to Z$ is a high resolution. Let $U_0\subset Z$ be the locus of $Z$ over which $Z'\to Z$ is an isomorphism.   Let $\Delta'$ be the sum of the strict transformations of $B$ and the reduced exceptional divisors of $X'\to X$ but with all components mapping outside $U_0$ removed. Let $U'_0\subset Z'$ be the preimage of $U_0$. Run an MMP on $K_{X'}+\Delta'$ over $Z'\times_Z X$ with some scaling, we end with a model $X''$ such that over $U'_0$ the pair $(X'',\Delta'')$ is a ${\mathbb Q}$-factorial dlt model of $(X,B)$ by the construction and $U'_0\times_{U_0}X|_{U_0}\simeq X|_{U_0}$. Hence $K_{X''}+\Delta''\sim_{\mathbb Q} f^*(K_X+B)\sim_{\mathbb Q} 0$ over $U_0'$ and $X''$ is of Fano type over $U_0'$. Run $K_{X''}+\Delta''$-MMP$/Z$ we end with a good minimal model over $Z'$ since every non-klt center of $(X'',\Delta'')$ is mapped into $U_0'$ and $X''$ is of Fano type over $U_0'$. We denote the final model again by $(X'',\Delta'')/Z'$ for convenience.
    
    Now we are going to show $qM_Z'$ is nef Cartier. We only need to show $qM_Z'$ is integral since $Z'$ is smooth. Let $f'':X''\to Z''/Z'$ be the Iitaka fibration of $K_{X''}+\Delta''$. On a common resolution $W$ of $(X,B)$ and $(X'',\Delta'')$. Then the pullback of $K_X+B$ and $K_{X''}+\Delta''$ will coincide over $U_0'$. Let $K_{X''}+B''$ and $L''$ be the pushdown from $W$ of the pullback of $K_X+B$ and $L$ respectively.  Let $U_0''\subset Z''$ be the preimage of $U_0'$, set $P'':=\Delta''- B''$ which is vertical and $\sim _{\mathbb Q} 0$ over $Z''$ since it is $0$ over $U_0''$, hence $P''=f''^*P_{Z''}$ for some ${\mathbb Q}$-Cartier ${\mathbb Q}$-divisor $P_{Z''}$. Denote $\Delta_{Z''}$ and $B_{Z''}$ to be the discriminant part of $(X'',\Delta'')/Z''$ and $(X'',B'')/Z''$ respectively, we see $\Delta_{Z''}=B_{Z''}+P_{Z''}$ by the construction of $P''$. Moreover  we have that \[\begin{aligned}q(K_{X''}+\Delta'')&=q(K_{X''}+B''+P'')\sim q(L''+P'')\\&=qf''^*(L_{Z''}+P_{Z''})= qf''^*(K_{Z''}+\Delta_{Z''}+M_{Z''})\end{aligned}.\] Here $L_{Z''}$ is the pullback of $L_Z$ to $Z''$ and $M_{Z''}$ is the moduli part of $(X'',\Delta'')/Z''$ and $(X'',B'')/Z''$ or equivalent saying $(X,B)/Z$.  The same arguments in the paragraphs above applied to $(X'',\Delta'')/Z''$, we see $qM_{Z''}$ is integral. So $qM_{Z'}$ is integral as well and we are done. 
    \end{proof}


Now we state an essential proposition which permits one to perturb the coefficients of the boundary which are close to $1$. The proof is the same as the characteristic $0$ case \cite[2.50]{birkar2019antipluricanonicalsystemsfanovarieties} since ACC for log canonical threshold in dimension $3$ and global ACC in dimension $2$ are known in characteristic $p >5$ \cite[1.10]{birkar2014existenceflipsminimalmodels}.
\begin{proposition}
\label{coe1}
    Let $\Phi\subset[0,1]$ be a DCC set, there is $\epsilon>0$ depending only on $\Phi$ satisfying the following. Let $(X,B)$ be a projective threefold pair with a contraction $X\to Z$ with $\dim Z>0$ such that:
    \begin{enumerate}
    \item $\chara k>5$
        \item $B\in\Phi\cup(1-\epsilon,1]$,
        \item $-(K_X+B)$ is a nef$/Z$ ${\mathbb Q}$-divisors,
        \item there is 
        \[0\leq P\sim_{\mathbb Q}-(K_X+B)/Z\]
        such that $(X, B+P)$ is generalized lc, and 
        \item $X$ is ${\mathbb Q}$-factorial and of Fano type over $Z$.
    \end{enumerate}
    Then let
    \[\Theta:=B^{\leq 1-\epsilon}+\lceil B^{>(1-\epsilon)}\rceil,\] run an $-(K_X+\Theta)$-MMP  over $Z$ and let $X'$ be the resulting model, then
    \begin{enumerate}
        \item $(X,\Theta)$ is lc, so is $(X',\Theta')$,
        \item the MMP does not contract any component of $\lfloor\Theta'\rfloor$,
        \item $-(K_{X'}+\Theta')$ is nef$/Z$.
    \end{enumerate}
\end{proposition}

Now we deal with the pl-contractions case which is what we want to reduce to.
\begin{lemma}
\label{plrcomp}
    Assume $(X,B)$ is a lc pair of threefolds and $f:X\to Z$ is a contraction with  $3>\dim Z>0$. Let $R$ be a finite set of rationals, then there is a natural number $n=n(R)$ such that suppose:
    \begin{enumerate}
        \item $\chara k>5$,
        \item $B\in R$,
        \item $X$ is of Fano type$/Z$,
        \item there is some ${\mathbb Q}$-divisor $\Gamma\in R$ such that $(X,\Gamma)$ is ${\mathbb Q}$-factorial plt,
        \item $S:=\lfloor \Gamma\rfloor$ is irreducible and a component of $\lfloor B\rfloor$ which is vertical$/Z$,
        \item $S$ intersect the fiber of $X/Z$ over some point $z\in Z$
        \item $-(K_X+\Gamma)$ is ample$/Z$,
        \item $-(K_X+B)$ is nef$/Z$, and
        \item $-S$ is nef$/Z$ if $Z$ is a curve.
    \end{enumerate}
    Then for such $z\in Z$, there is an $n$-complement $(X,B^+)$ of $(X,B)$ over $z$ with $B^+\geq B$.
\end{lemma}
\begin{proof}
    Since the statement is local, we may shrink near $f(S)$. If $-S$ is nef$/Z$, by \cite[1.1]{Bernasconiplvanishing} we see $R^if_*\mathcal O_X(-S)=0$ for $i>0$ since $-S-(K_X+\Gamma)$ is $f$-ample. If $Z$ is a surface but $-S$ is not necessarily nef$/Z$, then the same proof of \cite[2.2]{Bernasconiplvanishing} and \cite[2.8]{Bernasconibasepointfreeklt} apply similarly. That is, for any Weil divisor $L$ on $X$ such that $L-(K_X+\Gamma)$ is $f$-ample, we have $R^1f_*\mathcal O_X(L-S)=0$. So we set $L=\mathcal O_X$ here and again we see $R^1f_*\mathcal O_X(-S)=0$. We have the exact sequence 
    \[\mathcal O_Z=f_*\mathcal O_X\to f_*\mathcal O_S\to R^1f_*\mathcal O_X(-S),\] hence $\mathcal O_Z\to f_*\mathcal O_S$  is surjective. Let $\pi:V\to f(S)\to Z$ be the finite part of the Stein factorization of $S\to Z$, we have that:
    \[\mathcal O_Z\to\mathcal O_{f(S)}\to \pi_*\mathcal O_V=f_*\mathcal O_S\] is surjective, so $S\to f(S)$ is a contraction. 
    
    The aim of the following proof is to use the relative Kawamata-Viehweg vanishing to lift sections from $S$ to $X$. Suppose $\phi:X'\to X$ is a log resolution of $(X,B)$, let $S'$ be the birational transform of $S$ with the natural morphism $\psi:S'\to S$. Write \[K_S+B_S:=(K_X+B)|_S\] by divisorial adjunction, hence $(S,B_S)$ is klt and $B_S\in \Phi(S)$ for some finite set $S\subset [0,1]$ of rational numbers. $S$ is of Fano type, so by induction hypothesis, $(S,B_S)$ admits an $n$-complement $(S,B_S^+)$ over $z$ with $B_S^+>B_S$. Replacing $n$ by a multiple we may assume $nB$ and $n\Gamma$ are integral. We set \[N':=-(K_{X'}+B'):=-\phi^*(K_X+B),\quad S'\in T':=\lfloor B'^{\geq 0}\rfloor, \quad\Delta':=B'-T'.\] Define the integral divisor \[L':=-nK_{X'}-nT'-\lfloor(n+1)\Delta'\rfloor=n\Delta'-\lfloor(n+1)\Delta'\rfloor+nN'.\] Set \[K_{X'}+\Gamma':=\phi^*(K_X+\Gamma).\] Replace $\Gamma$ with $(1-a)\Gamma+aB$ for some $a<1$ sufficiently closed to $1$, we may assume $\Gamma'-B'$ has sufficiently small coefficients. Set \[P':=\sum\limits_{D'\text{ prime, }D'\neq S'}-\mu_{D'}(\lfloor\Gamma'+n\Delta'-\lfloor(n+1)\Delta'\rfloor\rfloor) D',\] to be a divisor on $X'$ and \[\Lambda':=\Gamma'+n\Delta'-\lfloor(n+1)\Delta'\rfloor+P'.\] 
    We claim that $P'\in\{0,1\}$, exceptional over $X$, and that $(X,\Lambda')$ is plt with $\lfloor \Lambda'\rfloor=S'$. In fact, let $D'$ be a component of $P'$, if it is a component of $T'$, then it is not a component of $\Delta'$ and hence $-\mu_{D'}(\lfloor\Gamma'+n\Delta'-\lfloor(n+1)\Delta'\rfloor\rfloor)=-\mu_{D'}(\lfloor \Gamma'\rfloor)=0$ as $D'\neq S'$. If it is not a component of $T$, then $\mu_{D'}(\lfloor\Gamma'+n\Delta'-\lfloor(n+1)\Delta'\rfloor\rfloor)=\mu_{D'}(\lfloor \Gamma'-\Delta'+\{(n+1)\Delta'\}\rfloor)=\mu_{D'}(\lfloor \Gamma'-B'+\{(n+1)\Delta'\}\rfloor)$, here $\{b\}$ means the fractional part of the real number $b$. Since $\Gamma'-B'$ has sufficiently small coefficients, we see $\Gamma'-B'+\{(n+1)\Delta'\}\in[-1,1)$, which means that $-\mu_{D'}(\lfloor \Gamma'-B'+\{(n+1)\Delta'\}\rfloor)\in\{0,1\}$ and hence $P'\in\{0,1\}$. Now we are going to show $P'$ is exceptional. In fact if $D'$ is a non exceptional component in $P'$, denote its image on $X$ to be $D$ then $\mu_{D'}P'=-\mu_D(\lfloor \Gamma-\Delta+\{(n+1)\Delta\})=-\mu_D(\lfloor \Gamma\rfloor)=0$ as $\Delta=\{B\}$ is effective and $n\Delta$ is integral and $D\neq S=\lfloor \Gamma\rfloor$.

    Let $R_S:=B_S^+-B_S$, then around $z$ we have \[-n(K_S+B_S)\sim-n(B_S^+-B_S)=nR_S\geq 0.\] Set $R_{S'}:=\psi^*R_S$, we see that \[-n(K_{S'}+B_{S'}):=-n\psi^*(K_S+B_S)\sim nR_{S'}\geq 0\] and \[nN'|_{S'}\sim nR_{S'}\] by pullback diagrams. Hence \[(L'+P')|_{S'}\sim G_{S'}:=n R_{S'}+n\Delta_{S'}-\lfloor (n+1)\Delta_{S'}\rfloor+P_{S'}.\] Moreover $G_{S'}\geq 0$ since no components $D'$ of $n\Delta'-\lfloor(n+1)\Delta'\rfloor$ has coefficient $\geq 1$ and $(X',B')$ is log smooth. 
    
    Now let $A':=\phi^*(-(K_X+\Gamma))$ which is pull back of an ample ${\mathbb Q}$-divisor$/Z$, we see \[L'+P'-S'=K_{X'}+\Gamma'+A'+n\Delta'-\lfloor(n+1)\Delta'\rfloor+nN'+P'-S'=K_{X'}+\Lambda'+A'+nN'-S'.\] We see that $A'+nN'$ is pull back of an ample ${\mathbb Q}$-divisor$/Z$ and 
    \[\phi_*(K_{X'}+\Lambda')=K_X+\Gamma+\phi_*(n\Delta'-\lfloor(n+1)\Delta'\rfloor)\leq K_X+\Gamma\] since $P$ is exceptional and $nB$ is integral,  by \Cref{van} we have $R^1f_*(L'+P'-S')=0$. Hence \[f_*(L'+P')\to (f|_S)_*((L'+P')|_{S'})\] is a surjection and hence $G_{S'}$ lifts to $G'$ over $z$ and support does not contain $S'$.  Pushing down $L',P',G',T',\Delta'$ to $X$, we get \[-nK_X-nT-\lfloor(n+1)\Delta\rfloor=L\sim G\geq 0.\] Since $nB$ is integral ,we have \[-n(K_X+B)=-nK_X-nT-n\Delta=L\sim nR:=G\geq 0,\] we set \[B^+:=B+R.\]
    Then $n(K_X+B^+)\sim 0$. Set \[nR':=G'-P'+\lfloor (n+1)\Delta'\rfloor-n\Delta'\sim nN'\sim_{\mathbb Q}0/X\] is the pull back of $nR$ which fills the pullback diagram for $R$. Hence \[K_S+B_S^+=(K_X+B^+)|_{S}.\] By inversion of adjunction, $(X,B^+)$ is lc near $S$. If $(X,B^+)$ is not lc near fiber over $z$, let \[\Omega:=aB^++(1-a)\Gamma\] for $a$ close to $1$ such that $(X,B^+)$ is not lc near fiber over $z$. Then $-(K_X+\Omega)$ is ample$/Z$ and the non-klt center near the fiber over $z$ has an extra component other than $S$ which contradicts the connectedness principle.
\end{proof}
Now we are going to prove the general case.
\begin{theorem}
\label{relacomp1}
    Assume $(X,B)$ is a 3-dimensional projective lc pair, $f:X\to Z$ is a contraction with $\dim Z>0$, let $R\subset [0,1]$ be a finite set of rational numbers, then there is some natural number $n=n(R)\in\mathbb N$ such that suppose:
    \begin{enumerate}
        \item $\chara k>5$
        \item $B\in\Phi(R)$,
        \item $X$ is of Fano type$/Z$,
        \item $-(K_X+B)$ is nef$/Z$, and is $\mathbb Q$-linearly trivial$/Z$ if $Z$ is a curve.
    \end{enumerate}
    Then, for any $z\in Z$, there is an $n$-complement $K_X+B^+$ of $K_X+B$ over $z$ with $B^+\geq B$.
\end{theorem}
\begin{proof}
    First, we pick an effective Cartier divisor $N$ on $Z$ passing through $z$ and set \[\Omega:=B+tf^*N,\] where \[t=\text{lct}_z(f^*N,X,B).\]  Since the statement is local, we shrink $Z$ so that $(X,\Omega)$ is lc everywhere. Take $(X',\Omega')$ to be a dlt ${\mathbb Q}$-factorial modification of $(X,\Omega)$. Then $X'$ is of Fano type and there is some $\Delta'\in\Phi(R)$ such that $\Delta'\leq\Omega'$ and some vertical component of $\lfloor\Delta'\rfloor$ intersects $X_z$ and $\Delta\geq B$ where $\Delta$ is the pushdown of $\Delta'$ to $X$. Run an MMP$/Z$ on \[-(K_{X'}+\Delta')=-(K_{X'}+\Omega')+(\Omega'-\Delta'),\] since $-(K_{X'}+\Omega')$ is nef$/Z$ and $\Omega'-\Delta'$ is effective, we end with a minimal model $(X'',\Delta'')$. It suffices to prove the boundedness of complements for $(X'',\Delta'')$, hence we may assume that $X$ is ${\mathbb Q}$-factorial and that some vertical component of $\lfloor B\rfloor$ intersects the fiber $X_z$.
    
    We now perturb the coefficients of $B$ so that $B\in R$ for some finite set $R$. Let $\epsilon>0$ be a sufficiently small number, let $\Theta$ be a boundary on $X$ defined by
    \[\Theta:=B^{\leq 1-\epsilon}+\lceil B^{>(1-\epsilon)}\rceil.\]
    Run MMP$/Z$ on $-(K_X+\Theta)$ and let $X'$ be the resulting model. By \Cref{coe1}, we can choose $\epsilon$ depending only on $R$ such that no component of $\lfloor\Theta\rfloor$ is contracted by the MMP, $(X',\Theta')$ is lc, and $-(K_{X'}+\Theta')$ is nef over $Z$. Moreover, the coefficients of $\Theta'$ belongs to some finite set $R$ since $1$ is the only accumulation point of $\Phi(R)$. If $K_{X'}+\Theta'$ has an $n$-complement over $z$, then so is $K_X+\Theta$ and so is $K_X+B$.  So by replacing $(X,B)$ with $(X',\Theta')$, we may assume that $B\in R$.  
    
    Fix an $\alpha<1$ ver close to $1$, set \[\Delta:=\sum
    \limits_{D\text{ is vertical over }Z}\mu_D(B)D+\sum\limits_{D\text{ is horizontal over }Z}\alpha\mu_D(B)D.\] Then $\Delta=\alpha B$ near the generic fiber and \[-(K_X+\Delta)=-\alpha(K_X+B)-(1-\alpha)K_X\] is big near the generic fiber, hence $-(K_X+\Delta)$ is big$/Z$. Let $g:X\to V/Z$ be the Iitaka fibration induced by $-(K_X+B)$, run $-(K_X+\Delta)$-MMP$/V$ we end with a minimal model $X'$. Then $-(K_{X'}+\Delta')$ is nef and big over $V$ but may not nef over $Z$. However replace $\Delta$ with $aB+(1-a)\Delta$ for $a$ sufficiently close to $1$, we may assume $-(K_{X'}+\Delta')$ is nef and big over $Z$. In fact, assume \[-(K_{X'}+\Delta')\sim_{\mathbb Q} g^*D+N'/Z\] for some nef and big ${\mathbb Q}$-divisor $N'/Z$ on $X'$ and some divisor $D$ on $V$, we have \[-(K_{X'}+B')\sim_{\mathbb Q}g^*A/Z\] for some ample ${\mathbb Q}$-divisor $A/Z$ on $V$ by definition, hence \[-(K_{X'}+aB'+(1-a)\Delta')\sim_{\mathbb Q}g^*(aA+(1-a)D)+(1-a)N'/Z.\] For $a$ very close to $1$, $-(K_{X'}+aB'+(1-a)\Delta')$ is nef and big over $Z$. 
    
    The MMP doesn't contract any component of $\lfloor \Delta\rfloor$. In fact if a component $S$ is contracted at some step, say $h:(X_1,\Delta_1)\to (X_2,\Delta_2)$, then $K_{X_1}+\Delta_1-h^*(K_{X_2}+\Delta_2)$ is $h$-nef and intersects the extremal ray positively, hence $-(K_{X_1}+\Delta_1)+h^*(K_{X_2}+\Delta_2)$ is effective by negativity lemma, moreover the coefficient of $S^\sim$ in $h^*(K_{X_2}+\Delta_2)$ is larger than $1$, which contradicts the fact that $(X,\Delta)$ is lc. By replacing $(X,B)$ with $(X',B')$ and $\Delta$ with $\Delta'$, we may assume that $-(K_X+\Delta)$ is nef and big over $Z$ 
    
    Let $\tilde\Delta:=\beta\Delta$ for some $\beta<1$ close to $1$, let $X\to T/Z$ be the contraction induced by $-(K_X+\Delta)$, replace $(X,B)$ by a dlt ${\mathbb Q}$-factorial modification and suppose $\alpha,\beta$ are sufficiently close to $1$, we run $-(K_X+\tilde\Delta)$-MMP$/T$ we end with some model $(X',B')$ with boundaries $\Delta',\tilde\Delta'$ such that $-(K_{X'}+\Delta')$ and $-(K_{X'}+\tilde\Delta')$ are nef and big$/Z$, $(X',B')$ is dlt ${\mathbb Q}$-factorial and $(X',\tilde\Delta')$ is klt, $\lfloor \Delta\rfloor$ is vertical and some component of it intersect $X'_z$. Replace $(X,B,\Delta,\tilde \Delta)$ with $(X',B',\Delta',\tilde \Delta')$ and shrink $Z$ around $z$ so that we may assume that every component of $\lfloor \Delta\rfloor$ intersects $X_z$.

    For further reduction, we may assume \[-(K_X+\Delta)\sim_{\mathbb Q}A+G/Z\] where $A$ is ample and $G$ is effective. If $\text{Supp}(G)$ contains no non-klt centers of $(X,\Delta)$, then $(X,\Delta+\delta G)$ is dlt for all sufficiently small $\delta>0$. Moreover \[-(K_X+\Delta+\delta G)\sim_{\mathbb Q}(1-\delta)(A+G)+\delta A/Z\] is ample$/Z$. Hence by perturbing the coefficients of $\Delta+\delta G$ we can make some $(X,\Gamma\sim_{\mathbb Q}\Delta+\delta G)$ such that $\lfloor\Gamma\rfloor=S$ is a vertical component of $\lfloor B\rfloor$ intersecting $X_z$ and $-(K_X+\Gamma)$ is ample$/Z$. If $Z$ is a surface, then apply \Cref{plrcomp} and we are done. Otherwise, we assume that $K_X+B\sim_\mathbb Q0/Z$ and that $Z$ is a curve. We run $-\upsilon S=K_X+\Gamma-\upsilon S-(K_X+\Gamma)$-MMP$/Z$ for some $\upsilon\ll 1$, we see $S$ is not contracted since the MMP is $S$-positive. Since $K_X+B\sim_{\mathbb Q}0/Z$, running $-\upsilon S$-MMP$/Z$ is $K_X+B$-crepant and  we only need to prove the boundedness of complements with the ending model by \Cref{subcrep}. Now we end with some minimal model $(X',B',\Gamma')$ with $-S'$-nef$/Z$ since $S$ is vertical, apply \Cref{plrcomp} and we are done.
    
    If $\text{Supp}(G)$ contains some non-klt center of $(X,\Delta)$, set \[\Omega:=\tilde\Delta+\text{lct}_z(G+\Delta-\tilde\Delta,X,\tilde\Delta) (G+\Delta-\tilde\Delta).\] By letting $\beta\to 1$ in the settings of $\tilde\Delta$ we may assume $t$ and $\Delta-\tilde\Delta$ are sufficiently small, we can assume every non-klt center of $(X,\Omega)$ is also a non-klt center of $(X,\Delta)$ since no extra centers will be added. Set $t:=\text{lct}_z(G+\Delta-\tilde\Delta,X,\tilde\Delta)$, we see that \[\begin{aligned}-(K_X+\Omega)&=-(K_X+\tilde\Delta+t(G+\Delta-\tilde\Delta))=-(K_X+\Delta)+\Delta-\tilde\Delta-t(G+\Delta-\tilde\Delta)\\&\sim_{{\mathbb Q},Z}A+G-tG+(1-t)(\Delta-\tilde\Delta)\\&=(1-t)(\frac{t}{1-t}A+A+G-((K_X+\Delta)-(K_X+\tilde\Delta))\\&\sim_{{\mathbb Q},Z}(1-t)(\frac{t}{1-t}A-(K_X+\tilde\Delta)),\end{aligned}\] which is ample over $Z$ since $-(K_X+\tilde\Delta)$ is nef and big over $Z$. 
    
    If $\lfloor \Omega\rfloor\neq 0$, then there is a vertical component $S$ of $\lfloor \Omega\rfloor\leq\lfloor\Delta\rfloor\leq\lfloor B\rfloor$, with the same discussion applied on $(X,\Omega)$ as the first case we are done. 
    
    If $\lfloor \Omega\rfloor=0$, let $(X',\Omega')$ be a dlt ${\mathbb Q}$-factorial modification, shrink $Z$ we assume all components of $\lfloor \Omega'\rfloor$ intersect $X_z$. Run MMP$/X$ on $K_{X'}+\lfloor\Omega'\rfloor$ we terminate at $X$ since $\lfloor \Omega'\rfloor$ are reduced exceptional divisor of $X'\to X$ and $X$ is ${\mathbb Q}$-factorial klt, if $X'\to X''/X$ ends with a minimal model different from $X$, then $X$ is lc but not klt, which is a contradiction. The last step is a divisorial contraction $X''\to X$ contracting some prime divisor $S''$ where $(X'',S'')$ is plt and $-(K_{X''}+S'')$ is ample$/X$. Denote $K_{X''}+\Omega''$ to be the pullback of $K_X+\Omega$, then $S''\in\lfloor\Omega''\rfloor$. Hence let $\Gamma'':=aS''+(1-a)\Omega''$ for $a>0$ sufficiently small, then $-(K_{X''}+\Gamma'')$ is a globally ample divisor$/Z$ and $(X,\Gamma'')$ is plt with $\lfloor \Gamma''\rfloor=S''$, running $-\upsilon S$-MMP if $Z$ is a curve and using \Cref{plrcomp} again and we are done.
\end{proof}

We now describe the progress on lifting complements from an lc trivial fibration structure $(X,B)\to V/Z$ with $\dim V\geq 1$ for Fano type threefolds over $Z$. Since \Cref{cplt} is solved completely for the case when the fibers are curves by \Cref{relacomp1}, we only need to consider the case when $Z$ is a curve or a point. We are going to use \Cref{effcbf} to show \Cref{cplt} for such $(X,B)/Z$. Firstly, we consider the case when $Z$ is a curve. If $V$ is also a curve, then $V\to Z$ is a contraction, hence $V=Z$. This case is solved by \Cref{relacomp1} already. So we may assume $V$ is a surface and we have the following result:

\begin{proposition}
\label{relacomp2}
    Let $R\subset [0,1]$ be a finite set of rational numbers, there is an $n\in\mathbb N$ and a prime number $p_0$ depending only on $R$ such that if 
    \begin{enumerate}
        \item $(X,B)$ is a projective threefold lc pair defined over an algebraically closed field $k$ with $\chara k=p>p_0$,
        \item $Z$ is a projective normal curve,
        \item there is a contraction $g:X\to Z$, $X$ is of Fano type over $Z$,
        \item $-(K_X+B)$ is nef over $Z$,
        \item there is a contraction $f:(X,B)\to V$ such that $(K_X+B)\sim_{\mathbb Q} 0/V$ with $\dim V=2$,
        \item $B\in\Phi(R)$.
    \end{enumerate} 
    Then there is an $n$-complement $(X,B^+)$ for $(X,B)$ with $B^+\geq B$ over arbitrary point $z\in Z$.
\end{proposition}
\begin{proof}
    By \Cref{effcbf}, we have the effective canonical bundle formula 
    \[q(K_X+B)\sim qf^*(K_V+B_V+M_V)\]
    where $B_V$ and $M_V$ are the discriminant part and moduli part of the canonical bundle formula applied to $(X,B)/V$ and $B_V\in\Phi(S)$ for some finite set $S$ of rationals in $[0,1]$ and $qM_{V'}$ is nef Cartier for any high model $V'\to V$.

    Now we are going to show that the generalized pair $(V,B_V+M_V)$  is generalized lc. Indeed take a high resolution $\psi:V'\to V$ and a resolution $\phi:(X',B')\to (X,B)$ where $\phi^*(K_X+B)=K_{X'}+B'$ such that $X'\to V'$ is a morphism and the moduli part $M_{V'}$ of $X'/V'$ is nef and $qM_{V'}$ is Cartier. Then we have:
    \[q(K_{X'}+B')\sim qf'^*(K_{V'}+B_{V'}+M_{V'})=qf'^*\psi^*(K_V+B_V+M_V)\]
    So it suffices to prove that $(V',B_{V'})$ is sub-lc. Since $(X,B)$ is lc, we have that $(X',B')$ is sub-lc. Thus for any prime divisor $D$, the coefficient of $D$ in $B_{V'}$ is \[1-lct_D(X',B',f^*D)\leq 1.\] So $(V',B_{V'})$ is sub-lc as desired. Moreover $V$ is of Fano type$/Z$ by \Cref{contr}. By induction hypothesis, $K_V+B_V+M_V$ has an $n$-complement $K_V+B^+_V+M_V$ over $z$ for some $n$ divisible by $q$ such that $G_V:=B^+_V-B_V\geq0$. Denote the pull back of $G_V$ to $V',X,X'$ by $G_{V'},G,G'$ respectively. Let $B^+:=B+G$, then $n(K_X+B^+)\sim 0$ over a neighbourhood of $z$ by definition of a complement. Thus we only need to prove that $(X,B^+)$ is lc.

    Let $C$ be a prime divisor over $X$. If $C$ is horizontal over $V$, then the statement holds  because $G$ is vertical. We therefore consider the case when $C$ is vertical over $V$ and defined on a high model $\phi:X'\to X$. Let $K_{X'}+B'+G'=\phi^*(K_X+B^+)$, we only need to show $\mu_C(K_{X'}+B'+G')\leq 1$. We prove this by contradiction. Assume there exists a vertical divisor $C$ such that $\mu_C(K_{X'}+B'+G')>1$. Since $(V,B_V^++M_V)$ is generalized lc, on the high model $V'$, we have
    \[K_{V'}+B_{V'}+G_{V'}+M_{V'}=\psi^*(K_V+B^+_V+M_V).\]
    Let $D$ be the image of $C$ on $V'$ since $C'$ is vertical, we have $\mu_D(B_{V'}+G_{V'})\leq 1$. Let $t_D:=\text{lct}_{\eta(D)}(X',B',f'^*D)$, we have that $\mu_D(G_{V'})\leq t_D$ since $\mu_D(B_{V'})=1-t_D$ by the constrcution of canonical bundle formula. Since $G'=f'^*G_{V'}$, $(X',B'+G')$ is sub-lc over $\eta(D)$, which leads to a contradiction.
\end{proof}
Now suppose $Z=\text{pt}$, that is we are going to construct a (global) bounded complements for Fano type threefold pairs $(X,B)$ which admits a lc trivial fibration structure $(X,B)\to V$. The following result is one of our main results and we have seen the similar proof in \Cref{relacomp2}.
\begin{theorem}
\label{fibcomp}
    Let $R\subset [0,1]$ be a finite set of rational numbers, there is an $n\in\mathbb N$ and a prime number $p_0$ depending only on $R$ such that if $(X,B)$ is a projective lc pair defined over an algebraically closed field $k$ satisfying the following conditions:
    \begin{enumerate}
        \item $\chara k=p>p_0$,
        \item $-(K_X+B)$ is nef,
        \item $X$ is of Fano type,
        \item $B\in \Phi(R)$ and
        \item there is a contraction $f:(X,B)\to V$ such that $(K_X+B)\sim_{\mathbb Q} 0/V$ with $3>\dim V>0$,
    \end{enumerate}
    then there is an $n$-complement $(X,B^+)$ for $(X,B)$ with $B^+\geq B$.
\end{theorem}
\begin{proof}
    We first consider the case when the fibers of $X/V$ are curves or $X/V$ is birational. In this case, by \Cref{effcbf}, we have the effective canonical bundle formula 
    \[q(K_X+B)\sim qf^*(K_V+B_V+M_V)\]
    where $B_V$ and $M_V$ are the discriminant part and moduli part of the canonical bundle formula applied to $(X,B)/V$ and $B_V\in\Phi( S)$ for some finite set $ S$ of rationals in $[0,1]$ and $qM_{V'}$ is nef Cartier for any high model $V'\to V$.

    Now we are going to show that the generalized pair $(V,B_V+M_V)$  is generalized lc. Indeed take a high resolution $\psi:V'\to V$ and a resolution $\phi:(X',B')\to (X,B)$ where $\phi^*(K_X+B)=K_{X'}+B'$ such that $X'\to V'$ is a morphism and the moduli part $M_{V'}$ of $X'/V'$ is nef and $qM_{V'}$ is Cartier. Then we have:
    \[q(K_{X'}+B')\sim qf'^*(K_{V'}+B_{V'}+M_{V'})=qf'^*\psi^*(K_V+B_V+M_V)\]
    So it suffices to prove that $(V',B_{V'})$ is sub-lc. Since $(X,B)$ is lc, we have that $(X',B')$ is sub-lc. Thus for any prime divisor $D$, the coefficient of $D$ in $B_{V'}$ is \[1-lct_D(X',B',f^*D)\leq 1.\] So $(V',B_{V'})$ is sub-lc as desired. Moreover $V$ is of Fano type by \Cref{contr}. By induction hypothesis, $K_V+B_V+M_V$ has an $n$-complement $K_V+B^+_V+M_V$ for some $n$ divisible by $q$ such that $G_V:=B^+_V-B_V\geq0$. Denote the pull back of $G_V$ to $V',X,X'$ by $G_{V'},G,G'$ respectively. Let $B^+:=B+G$, then $n(K_X+B^+)\sim 0$ by definition of a complement. Thus we only need to prove that $(X,B^+)$ is lc.

    Let $C$ be a prime divisor over $X$. If $C$ is horizontal, the statement holds because $G$ is vertical. We therefore consider the case when $C$ is vertical and defined on a high model $\phi:X'\to X$. Let $K_{X'}+B'+G'=\phi^*(K_X+B^+)$, we only need to show $\mu_C(K_{X'}+B'+G')\leq 1$. We prove this by contradiction. Assume that there exists a vertical divisor $C$ such that $\mu_C(K_{X'}+B'+G')>1$. Since $(V,B_V^++M_V)$ is generalized lc, on the high model $V'$, we have
    \[K_{V'}+B_{V'}+G_{V'}+M_{V'}=\psi^*(K_V+B^+_V+M_V).\]
    Let $D$ be the image of $C$ on $V'$ since $C'$ is vertical, we have $\mu_D(B_{V'}+G_{V'})\leq 1$. Let \[t_D:=\text{lct}_{\eta(D)}(X',B',f'^*D),\] we have that $\mu_D(G_{V'})\leq t_D$ since $\mu_D(B_{V'})=1-t_D$ by the constrcution of canonical bundle formula. Since $G'=f'^*G_{V'}$, $(X',B'+G')$ is sub-lc over $\eta(D)$, a contradiction.

    Now consider the case when the fiber of $(X,B)/V$ are surfaces. Take a dlt modification of $(X,B)$ we may assume $(X,B)$ is dlt and ${\mathbb Q}$-factorial. By \Cref{singgen}, we have that the set of discrepancies near the generic fiber $A_\eta(X/Z,B)=\{0,>\epsilon\text{ part}\}$. Extract the divisors $D$ such that $a_\eta(D,X,B)\leq \epsilon$ we get $X'\to X$. Here $X'$ is of Fano type and $\epsilon$-lc. Run an $K_{X'}$-MMP$/V$, we get a Mori fiber space structure $X''\to Z/V$. If $\dim Z=2$, then $K_{X''}+B''\sim_{\mathbb Q}0/V$ and it falls in the first case. Otherwise, $X''\to Z$ is a Mori fiber space since $Z$ is a curve and $X_\eta''$ is $\epsilon$-lc. If $(X'',B'')/Z$ is not the exceptional case in \Cref{contr}, apply \Cref{contr} and the same proof above we are done. If $(X'',B'')/Z$ is the only exceptional case in \Cref{contr}, then we see that $K_Z\sim0$ is a $1$-complement already, enlarge $n$ such that $q$ divides $n$ and the same proof will also apply.
\end{proof}
As a corollary, we can prove our another main result:
\begin{corollary}
\label{fibcomp2}
    

    Let $R\subset [0,1]$ be a finite set of rational numbers. Then there is a constant $n$ and a prime number $p_0$ depending only on $R$, such that if $(X,B)$ is a $3$-dimensional projective lc pair defined over an algebraically closed field $k$ satisfying the following conditions:
    \begin{enumerate}
        \item $\chara k=p>p_0$,
        \item $B\in \Phi(R)$,
        \item $X$ is Fano type and
        \item $-(K_X+B)\not\equiv0$ is nef but not big,
    \end{enumerate}
    then there is an $n$-complement $(X,B^+)$ with $B^+\geq B$.
\end{corollary}
\begin{proof}
    The result follows by considering the fibration defined by $-(K_X+B)$ and applying \Cref{fibcomp}.
\end{proof}
\begin{remark}
    Assume good theory of generalized pairs for threefolds in positive characteristic, the author believes that the arguments above could apply to the generalized pair case $(X,B+M)\to V$ when the moduli part $M$ is not big$/V$.
\end{remark}

\section{Boundedness of relative complements}









%%%%%%%%%%%%%%%%%%%%%%%%%%%%%%%%%%%%%%%%%%%%%%%%%%%%%%%%%%%%%%%%%%%%%%%%%%%%%%%%%%%
To prove the boundedness of relative complements of a 3-dimensional Fano type fibration $f:(X,B)\to Z$, the only remaining case is when $Z$ is a curve and $-(K_X+B)$ is nef and big over $Z$. Taking the semiample fibration defined by $-(K_X+B)$, we may assume that $-(K_X+B)$ is ample over $Z$. Moreover, since the problem is local near $z\in Z$, we may shrink $Z$ and assume $Z$ is affine. We use the technique of Frobenius stable sections to lift complements, following a similar approach and setup as in \cite{threedimmmp} and \cite{HW19}. Let $f:X\to Y$ be a projective morphism between normal varieties over $k$, we denote the relative Frobenius stable section of a line bundle $L$ with respect to the Frobenius trace map associated to $(X,\Delta)/Y$ to be $S^0f_*(X,\Delta;L):=S^0f_*(\sigma(X,\Delta)\otimes \mathcal O_X(L))$, see \cite{threedimmmp} for the detailed definition.

Let $(X,S+B)/Z$ be a dlt log Fano fibration, here $S$ is reduced (might be empty) and $B<1$. Fix a natural number $m\in\mathbb N$ and we denote $A:=-(K_X+S+B)$, we study the sections in $f_*(\mathcal O_X(\lfloor mA\rfloor))$. If $B\in \Phi(R)$ for some finite set of rational numbers $R$, let $I$ be the least common denominator of $R$ and we may assume $I$ divides $m$. We define:
\[\begin{aligned}
    \Phi&:=S+\{(m+1)B\},\\
    L&:=\lfloor mA\rfloor.
\end{aligned}\]
For future perturbations, we choose an effective $\mathbb Q$-divisor $\Lambda$ with sufficiently small coefficients, having no common components with $S$, and such that the Cartier index of $K_X+S+B+\Lambda^m$ is not divisible by $p=\chara k$, where $\Lambda^m:=\frac{1}{m+1}\Lambda$. Since $m$ is divisible by $I(R)$ and $B\in\Phi(R)$, we see the defect divisor \[D:=\lceil mB\rceil-\lfloor (m+1)B\rfloor=0.\]
Hence we have 
\[\begin{aligned}
    \lfloor mA\rfloor&=-m(K_X+S)-\lceil mB\rceil\\
    &=-m(K_X+S)-\lfloor(m+1)B\rfloor\\
    &=K_X+\Phi-(m+1)(K_X+S+B,)
\end{aligned}\]
so $L-(K_X+\Phi+\Lambda)=-(m+1)(K_X+S+B+\Lambda^m)$ is ample$/Z$ with Cartier index non-divisible by $p$, and so the denominators appearing in $K_X+\Phi+\Lambda$ is not divisible by $p$.
\begin{definition}
    With the notations above, we define $C^0_\Lambda f_*(X,S+B;L):=S^0f_*(X,\Phi+\Lambda; L)\subset f_*(\mathcal O_X(L))$ the relative perturbed Frobenius stable section of $L$ associated to $(X,S+B)$.
\end{definition}

Consider a coefficient $a=1-\frac{r}{ln}\in\Phi(R)$ for some $\frac{r}{l}\in R$, may assume $l=m$ since $m$ is divisible by $l$, we see \[\{(m+1)a\}=\{m-\frac{r}{n}+1-\frac{r}{mn}\}=\{-\frac{r}{n}-\frac{r}{mn}\}\leq 1-\frac{r}{mn}=a.\] As a result, we see that $B\geq \Phi$. Let $T$ be the normalization of an irreducible component $T'\subset S$, define $K_T+B_T\sim_\mathbb Q(K_X+S+B)|_T$ and $K_T+\Phi_T+\Lambda_T\sim_\mathbb Q (K_X+\Phi+\Lambda)|_T$ be the divisorial adjunction formulas. Since any component of $S$ is not contained in $\Supp(\Lambda)$ and $\Lambda$ is sufficiently small, $B\geq \Phi$ implies that if $(T,B_T)$ is globally F-regular (resp. globally F-split), then $(T,\Phi_T+\Lambda_T)$ is also globally F-regular (resp. globally F-split).

Now recall our settings, $(X,B)\to Z$ is a 3-dimensional log Fano fibration over an affine curve. Take a $\mathbb Q$-factorial dlt modification we may assume $(X,S+B)$ is dlt $\mathbb Q$-factorial with $\lfloor S+B\rfloor=S$. We first consider the case when $S\neq0$ is irreducible and is vertical over $Z$ and $(X,S+B)$ is plt. In particular, $S$ is normal. Consider the adjunction formula $K_S+B_S:=(K_X+S+B)|_S$, we have that $(S,B_S)$ is a klt log del pezzo pair. We first prove the global F-regularity of $(S,B_S)$ when it is $\epsilon$-lc.

\begin{lemma}\label{sfrgfr}
    Let $X\to T$ be an affine flat morphism of normal varieties over an F-finite field of characteristic $p>0$, then $X$ is strongly F-regular if and only if $X$ is globally F-regular over $T$.
\end{lemma}
\begin{proof}
    Since relative globally F-regularity is a local condition, we may assume $T$ is affine, hence $X$ is affine and the assertion follows directly.
\end{proof}
\begin{definition}[Relative section ring]
    Let $f:X\to T$ be a contraction of normal varieties over an algebraically closed field, then for any line bundle $L$ on $X$ which is ample over $T$, we define the section $\mathcal O_T$-algebra, or the relative section ring associated to $(X,L)/T$ to be the following graded $\mathcal O_T$-module:
    \[\mathcal S:=\bigoplus\limits_{n\in \mathbb N}R^0f_*\mathcal O_X(nL).\]
    Here the $\mathcal O_T$-algebra structure on $\mathcal S$ is given naturally by the following morphism:
    \[R^0f_*\mathcal O_X(nL)\otimes R^0f_*\mathcal O_X(mL)\to R^0f_*(\mathcal O_X(nL)\otimes\mathcal O_X(mL))\simeq R^0f_*(\mathcal O_X((n+m)L)).\]
    We define the distinguished graded sub-module of $\mathcal S$ to be
    \[\mathcal S_+:=\bigoplus\limits_{n\in \mathbb N^+}R^0f_*\mathcal O_X(nL).\]
    This is a homogeneous ideal module of $\mathcal S$ and we have $\mathcal S/\mathcal S_+\simeq \mathcal O_T$ since $X/T$ is a contraction. We define the punctured spectrum of $\mathcal S$ to be $\Spec_{\mathcal O_T}\mathcal S-V(\mathcal S_+)$, here $V(\mathcal S_+)$ is the set consisting of all the prime ideals containing $\mathcal S^+$ after base changing to any affine open subscheme of $T$. 
\end{definition}
\begin{lemma}
    $V(\mathcal S_+)$ is a well-defined closed subset of $\Spec_{\mathcal O_T}\mathcal S$. Moreover, every ideal in $V(\mathcal S_+)$ is a homogeneous ideal.
\end{lemma}
\begin{proof}
    Let $\coprod\limits T_i\to T$ be an open affine covering of $T$ and $\coprod\limits_{k}T_{ijk}\to T_{ij}:=T_i\cap T_j$ to be an open affine covering of $T_{i,j}$. We only need to show $V(\mathcal S_+)$ is compatible with the hypercoverings. We may assume $T_i=\Spec A$ and $T_{ijk}:=\Spec A_f$ to be a principal open subscheme. Let $\mathcal S|_{T_i}=:S$ to be a graded $A$-algebra, and similar for $S_+$. Then $(V(\mathcal S_+)|_{T_i})\times_{T_i}T_{ijk}=\{\text{prime ideals in } S/S_+\simeq A\text{ with no intersection with }M(f)\}=\{\text{prime ideals in } A_f\}=V(\mathcal S_+)|_{T_{ijk}}$ is defined only depending on $T_{ijk}$, where $M(f):=\{1,f,f^2,\cdots\}$. Hence $V(\mathcal S_+)$ is compatible with the hypercoverings. The second argument is clear since $\mathcal S_+$ contains all the terms with grading larger than $0$.
\end{proof}
\begin{lemma}\label{sectionring}
    There is a naturally defined morphism $\Spec_{\mathcal O_T}\mathcal S-V(\mathcal S_+)\to X/T$, which can be viewed as a $\mathbb G_m(T):=\Spec_{\mathcal O_T}\mathcal O_T[t,t^{-1}]$ bundle over $X$ and there is a natural identification $X\simeq \Proj_T\mathcal S$.
\end{lemma}
\begin{proof}
    The identification is proved in \cite[II.4.6.2]{PMIHES_1961__8__5_0}. The first morphism is defined naturally by the identification.
\end{proof}
\begin{remark}
    The relative section algebra is compatible with the fibers. In fact, if $X/T$ has normal geometric fibers and $X_f$ is a general fiber of $X/T$, then the base change of the data $(X,L,\mathcal S,\mathcal S_+,\Spec_{\mathcal O_T}\mathcal S-V(\mathcal S_+)\to X=\Proj_T \mathcal S)$ to $X_f$ is just the global section ring construction $(X_f,L|_{X_f},S,\mathfrak m,\Spec S-\mathfrak m\to X_f=\Proj S)$. One can view $\Spec_{\mathcal O_T}\mathcal S/T$ to be the relative affine cone of $X$ in some sense.
\end{remark}
\begin{lemma}\label{étaleness}
      Let $\mathcal S^{(n)}$ denote the relative section ring of $X$ associated to $nL$ such that $\mathcal S^{(n)}$ is generated by degree-1 terms, we have a natural morphism $\Spec_{\mathcal O_T}\mathcal S\to \Spec_{\mathcal O_T}\mathcal S^{(n)}$. Suppose $n$ is not divisible by $p=\chara k$, then this induces an étale map on the punctured spectrum.
\end{lemma}
\begin{proof}
    Since the statement is local, we may assume $T$ is affine.   We can choose an open affine cover of $X$ which trivializes $L$ over $T$. That is, thinking of $X=\Proj_T(\mathcal S^{(n)})$, we may choose a homogeneous element $x\in \mathcal S^{(n)}$ and assume the covering has the form $D_+(x)=\Spec A_0$, where $A_0=[\mathcal S^{(n)}_x]_0$. Let $y\in X$ be a point in $D_+(x)$ we are looking at, suppose $L$ is trivial on an affine neighbourhood $U$ of $y$, then so it is on $U\cap D_+(x)$. Then there is some $f\in A_0$ such that $D_+(x,f):=\Spec (A_0)_f\subset U\cap D_+(x)$.  By definition, $f=\frac{g}{x^m}$, where $g$ is a homogeneous element in $\mathcal S$ of grading $mn$, so we have $D_+(xg)\subset U\cap D_+(x)$. Replace $x$ with $xg$, we can always find some $D_+(x)$ such that $L$ is trivial on $D_+(x)$.  Now set $D(x)=\Spec A=\Spec \mathcal S^{(n)}_x\to D_+(x)$ be the projection map, we see that $D(x)$ covers the punctured spectrum and hence we only need to prove that the morphism
    \[\mathcal S^{(n)}_x\to \mathcal S_x\] is étale. Since $L=\mathcal S(1)^\sim$, we have that \[L^n|_{D_+(x)}\simeq \mathcal S(n)[\frac{1}{x}]^{\sim}\simeq [\mathcal S[\frac{1}{x}]]_n\simeq [\mathcal S[\frac{1}{x}]]_{0}\] for any $n\in\mathbb Z$, where the last isomorphism follows by the triviality of $L$ on $D_+(x)$. Hence take a generator $t\in \mathcal S_1$, we see $\mathcal S_x\simeq \mathcal S^{(n)}_{(x)}[t,t^{-1}]$ and $\mathcal S^{(n)}_x=\mathcal S^{(n)}_{(x)}[t^n,t^{-n}]$ and the assertion follows since $n$ is not divisible by $p$.
\end{proof}
For any prime divisor $D$ on $X=\Proj_T \mathcal S$, it corresponds to a homogeneous prime ideal of $\mathcal S$ of height $1$, which induces a prime divisor on $\Spec_T \mathcal S$, we denote it by $D_\mathcal S$. Under this correspondence, we have $(K_{X})_\mathcal S=K_{\Spec_T\mathcal S}$. 
\begin{lemma}\label{kltcorresp}
    Suppose $(X,D)/T$ is a flat non-birational contraction, $\mathcal L$ is an ample$/T$ line bundle on $X$, $(\Spec_T\mathcal S,D_\mathcal S)$ and $(\Spec_T\mathcal S^{(n)},D_{\mathcal S^{(n)}})$ are the corresponding pairs defined by $(X,D)$, then $(\Spec_T\mathcal S,D_\mathcal S)$ is klt (resp. lc) if and only if $(\Spec_T\mathcal S^{(n)},D_{\mathcal S^{(n)}})$ is klt (resp. lc).
\end{lemma}
\begin{proof}
    Consider the morphism $f:\Spec_{T}\mathcal S\to \Spec_{T}\mathcal S^{(n)}$, we see that it is étale outside $V(\mathcal S_+)$ by \Cref{étaleness}. Since $X/T$ is flat and non-birational, $V(\mathcal S_+)$ is of codimension greater or equal than two in $\Spec_{T}\mathcal S$. So the finite morphism $f$ is étale in codimension $1$. Moreover, we see $f^*K_{\Spec_T\mathcal S^{(n)}}=K_{\Spec_T\mathcal S}$ and $f^*D_{\Spec_T\mathcal S^{(n)}}=D_{\Spec_T\mathcal S}$, so the log discrepancies will retain.
\end{proof}
\begin{theorem}\label{corresp}
    Let $(X,D)/T$ be a flat contraction over a field of characteristic $0$ and $T$ is smooth with $(X,D)$ klt, $(\Spec_T \mathcal S,D_S)$ is one of its section ring spectrum, then $(X,D)/T$ is a Fano fibration if and only if there exists a section ring $\mathcal S=\mathcal S(X,L)$ such that $(\Spec_T \mathcal S,D_S)$ is klt.
\end{theorem}
\begin{proof}
    Since the statement is local on $T$, we can shrink $T$ arbitrarily. By \Cref{kltcorresp}, we may assume $\mathcal S$ is generated by grading-1 terms. We denote $S$ to be $\Spec_T\mathcal S$ and $S^*$ to be $\Spec_T\mathcal S-V(\mathcal S_+)$  for convenience. We also denote $Y$ to be the blow up of $S$ along $V(\mathcal S_+)$. Since $T$ is smooth, this is a blow up along a smooth locus with the exceptional locus $X$. $Y\to X$ can be viewed as the tautological bundle of $X=\Proj_T \mathcal S$. Since $T$ is smooth, $S^*\to X$ is a smooth morphism, so $(X,D)$ is klt if and only if $(S^*,D_{S^*})$ is klt. 

    Suppose $K_S+D_S$ is $\mathbb Q$-Cartier on $S$, we assume $n(K_S+D_S)$ is Cartier, which means this is an $\mathcal S$-module locally free of rank $1$. We pick a point $z\in T$ and consider the problem locally near $z$. $n(K_S+D_S)$ is graded locally free means that near $z$, there is a open covering of principal open sub-schemes $\{D(x_i)\to S\}$, such that \[n(K_S+D_S)|_{D(x_i)}=\mathcal S_{(x_i)}(t_i)\] for some integer $t_i$ and homogeneous $x_i$ and $\{D(x_i)\to S\}$ covers $S_z$. Consider the restriction of $n(K_S+D_S)$ to $D(x_ix_j)\subset D(x_i),D(x_j)$, we see $t_i=t_j$. Hence all $t_i=t$, and $n(K_S+D_S)=\mathcal S(t)$ near $S_z$. So one can always shrink $T$ to assume $(K_S+D_S)|_{V(\mathcal S_+)}$ is $\mathbb Q$-linearly trivial and make $n(K_S+D_S)$ a free $\mathcal S$-module $\mathcal S(t)$. Then we have $K_S+D_S\sim_\mathbb Q\frac{t}{n}\mathcal O(1)$ is either ample, numerical trivial or anti-ample. Consider the map $\pi:Y\to S$, we have that \[\pi^*(K_S+D_S)\sim_\mathbb QK_Y+D_Y+aX,\] where $X$ is the exceptional divisor. We see that $X|_X=\mathcal O(-1)$ since $Y$ is the tautological bundle of $X$ with respect to $L=\mathcal S(1)$ defining $\mathcal S$ on $X=\Proj_T\mathcal S$, where $X\to Y$ is the zero section. Hence \[\begin{aligned}(K_S+D_S)|_X&\sim_\mathbb Q (K_Y+X)|_X+D_Y|_X+(a-1)X|_X\\&\sim_\mathbb Q K_X+D_X+(1-a)L\sim_\mathbb Q\frac{t}{n}L+(1-a)L/T.\end{aligned}\] On the other hand, \[(K_S+D_S)|_X=((K_S+D_S)|_{V(\mathcal S_+)})|_X\sim_\mathbb Q0/T. \] Hence we have $a=1+\frac{t}{n}$. Under this reduction, we see if $(S,D_S)$ is klt, then the above arguments apply and $a<1$ implies $\frac{t}{n}<0$, hence $K_X+D\sim_\mathbb Q\frac{t}{n}L$ is anti-ample over $T$. Conversely, if $(X,D)$ is log Fano$/T$, then $L:=-n(K_X+D)$ is an ample Cartier divisor$/T$ for some $n\in\mathbb N$, so the section ring associated to $L$ will give what we need.
\end{proof}
\begin{theorem}\label{affgfr}
    Let $X/T$ be a flat contraction of normal quasi-projective varieties over an algebraically closed field of characteristic $p$, $\mathcal S=\mathcal S(X,L)$ is a section ring associated to any ample$/T$ line bundle $L$ on $X$. Then the pair $(X,D)$ is globally F-regular over $T$ if and only if $(\Spec_T\mathcal S,D_S)$ is globally F-regular over $T$.
\end{theorem}
\begin{proof}
    The statement is local on $T$, so the arguments of \cite[5.3]{SS09} apply after shrinking $T$ to a sufficiently small affine open subschemes.
\end{proof}
Combining the above preparations, we obtain the following result.
\begin{theorem}\label{relgfr}
    Let $(X,\Delta)\to T$ be a flat contraction of normal varieties in characteristic $0$, suppose $(X,\Delta)$ is log Fano over $T$, then for its any integral model $(\mathfrak X,\daleth)/\mathfrak T$, $(\mathfrak X_p,\daleth_p)$ is globally F-regular over $\mathfrak T_p$ for almost all $p$. 
\end{theorem}
\begin{proof}
    By \Cref{corresp}, there exists a section ring spectrum $(\Spec_T\mathcal S,\Delta_S)/T$ which is klt. Hence by \cite[3.4]{tak04}, for almost all mod $p$ reductions, $(\Spec_T\mathcal S,\Delta_S)/T$ is strongly F-regular, hence globally F-regular$/\mathfrak T_p$ by \Cref{sfrgfr}, hence $(\mathfrak X_p,\daleth_p)$ is re-constructed to be globally F-regular$/\mathfrak T_p$ for almost all $p$ by \Cref{affgfr}.
\end{proof}
\begin{corollary}\label{gengfr}
    Let $\epsilon>0$ be a real number and $I\subset [0,1]\cap \mathbb Q$ be a finite set of rational numbers, then there exists a $p_0=p_0(\epsilon,I)$ such that for all $p>p_0$ and $\chara k=p$, an $\epsilon$-lc log Fano surface pair $(X,B)/k$ with $B$ in $I$ is globally F-regular.
\end{corollary}
\begin{proof}
    We assume $p_0>5$. Since $\epsilon$-lc log del Pezzo surfaces form a bounded family over $\mathbb Z[\frac{1}{30}]$, hence there is a flat projective family $(\mathfrak X,\mathfrak B)\to \mathfrak T$ of finite type over $\mathbb Z[\frac{1}{30}]$ such that each $(X,B)$ is one of its geometric fiber. Consider the locus $\mathfrak T'$ inside $\mathfrak T$ where $(\mathfrak X,\mathfrak B)$ is log Fano over $\mathfrak T$, which is an open subset of $\mathfrak T$, hence we can replace $\mathfrak T$ by $\mathfrak T'$. We see that $(\mathfrak X_0,\mathfrak B_0)$ is log Fano over $\mathfrak T_0$, hence $(\mathfrak X_p,\mathfrak B_p)$ is globally F-regular over $\mathfrak T_p$ for almost all $p$ by \Cref{relgfr}. Now taking geometric fibers and we are done.
\end{proof}
Now let us come back to our problem on the boundedness of complements.
\begin{definition}
    Let $(X,B)$ be a lc pair, we say it is $(0,\epsilon)$-lc (resp. near $S\subset \lfloor B\rfloor$) if it is not klt and its second minimal log discrepancy (resp. near $S$) is larger or equal than $\epsilon$.
\end{definition}
\begin{lemma}
\label{relacomp3}
     Assume $(X,S+B)$ is a 3-dimensional dlt pair, $f:X\to Z$ is a contraction. Let $\epsilon>0$ be a real number and $I\subset[0,1]$ be a finite set of rational numbers, then there exists a natural number $n$ and a prime number $p_0$ depending only on $\epsilon$ and $I$, such that suppose:
     \begin{enumerate}
        \item $\chara k=p>p_0$,
         \item $(X,S+B)$ is weak Fano over $Z$,
         \item $S+B\in \Phi(I)$,
         \item $\dim Z\leq 1$,
         \item $S$ is an irreducible reduced divisor, moreover $S$ is vertical over $Z$ if $Z$ is a curve, and
         \item $(X,S+B)$ is $(0,\epsilon)$-lc near $S$ and plt near $S$, moreover $\lfloor S+B\rfloor= S$ near $S$,
     \end{enumerate}  
     Then there exists an $n$-complement $(X,S+B^+)$ of $(X,S+B)$ over $z=f(S)$ with $B^+\geq B$.
\end{lemma}
\begin{proof}
    By \cite[2.18]{threedimmmp}, we see the restriction map $S^0f_*(X,\Psi+\Lambda;L)\to S^0(S,\Psi_S+\Lambda_S; L_S)$ is a surjection. Since $\Lambda_S=0$ and $\Psi_S=B_S$ and $(S,B_S)$ is $\epsilon$-lc weak del pezzo surface with $B_S\in J$, by \Cref{gengfr}, $(S,B_S)$ is globally F-regular. In particular $S^0(S,\Psi_S+\Lambda_S; L_S)=H^0(S,L_S)$. Enlarging $m$ so that $(S,B_S)$ admits an $m$-complement $(S, B_S+G_S)$ where $mG_S\sim L_S$. By surjectivity, we see that there is some $\Delta\in |L|$ such that $\Delta|_S=mG_S$. Define $G:=\frac{1}{m}\Delta$, we claim that $(X,S+B+G)$ is an $m$-complement of $(X,S+B)$. Indeed we see that $(X,S+B+G)$ is lc near $S$ by inversion of adjunction.  Suppose it is not lc near $z$, then we see that the lc center does not intersect $S$, hence for sufficiently small $\delta\geq 0$, we see that $(X,S+B+\delta G)$ is lc. So let $\lambda=\text{lct}(X,S+B,G)$, then $(X,S+B+\lambda G)$ is lc with two disjoint lc center mapped onto $z$ and $-(K_X+S+B+\lambda G)$ is nef and big over a neighbourhood of $z$, which contradicts the connectedness principle and we are done.
\end{proof}
Despite that we have not fully resolved with the case when $(X,S+B)$ is plt, we shall use the same method above to resolve the non-plt case first.
\begin{lemma}
\label{relacomp4}
     Assume $(X,S+B)$ is a 3-dimensional dlt pair, $f:X\to Z$ is a contraction. Let $I\subset[0,1]$ be a finite set of rational numbers, then there exists a natural number $n$ and a prime number $p_0$ depending only on $I$ such that suppose:
     \begin{enumerate}
         \item $\chara k=p>p_0$,
         \item $(X,S+B)$ is weak Fano over $Z$,
         \item $S+B\in \Phi(I)$,
         \item $\dim Z\leq 1$,
         \item near $S$, a lc center $T\subset \lfloor S +B\rfloor=S$ is an irreducible reduced divisor, moreover $T$ is vertical if $Z$ is a curve, and
         \item $(X,S+B)$ is not plt near $S$.
     \end{enumerate}
     Then there exists an $n$-complement $(X,S+B^+)$ of $(X,S+B)$ over $z=f(T)$ with $B^+\geq B$.
\end{lemma}
\begin{proof}
    According to the above proof, replacing $T$ with its normalization, it suffices to prove that $(T,\Phi_T+\Lambda_T)$ is globally F-split. However, this follows by global F-adjunction applied to a non-klt center of $(T,\Phi_T+\Lambda_T)$ immediately since $(T,\Phi_T+\Lambda_T)$ is a non-klt del Pezzo type log surface pair by divisorial adjunction, cf. \cite[Subsection 2.2]{CTW16}.
\end{proof}

\begin{lemma}
\label{relacomp5}
     Assume $(X,S+B)$ is a 3-dimensional dlt pair, $f:X\to Z$ is a contraction. Let  $I\subset[0,1]$ be a finite set of rational numbers, then there exists a natural number $n$, a small real number $\epsilon>0$ and a prime number $p_0$ depending only on $I$ such that suppose:
     \begin{enumerate}
         \item $\chara k=p>p_0$,
         \item $(X,S+B)$ is weak Fano over $Z$,
         \item $S+B\in \Phi(I)$,
         \item $\dim Z= 1$,
         \item $S$ is an irreducible reduced vertical divisor over $Z$,
         \item $(X,S+B)$ is plt near $S$, moreover $\lfloor S+B\rfloor= S$ near $S$, and
         \item $(X,S+B)$ is not $(0,\epsilon)$-lc near $S$.
     \end{enumerate}
      Then there exists an $n$-complement $(X,S+B^+)$ of $(X,S+B)$ over $z=f(S)$ with $B^+\geq B$.
\end{lemma}
\begin{proof}
    We may assume that $(X,S+B)$ has a place whose log discrepancy is in $(0,\epsilon)$ near $S$, say $D$. Let $X'\to X$ be a birational contraction which extracts $D$, we see $f^*(K_X+S+B)=K_{X'}+S'+B'+bD'$ with $b>1-\epsilon$ and $D$ intersects $S$. By ACC of lct, we see $\lct(X',S'+B', D')$ belongs to an ACC set, hence $(X',B'+S'+D')$ is lc so we replace $X$ with $X'$. We may assume $(X,S+B+D)$ is dlt by taking a dlt $\mathbb Q$-factorialization. Now we run $-(K_X+B+S+D)$-MMP$/Z$ on $X$, we see that $D$ is not contracted since the MMP is $D$-positive. Now we claim that $S$ will not be contracted under this MMP if $\epsilon$ is small enough. In fact, suppose at some step $g:X_i\to X_{i+1}$, $S$ is contracted. We see $(K_{X_i}+S_i+B_i+D_i)-g^*(K_{X_{i+1}}+S_{i+1}+B_{i+1}+D_{i+1})$ is $g$-ample, hence $g^*(K_{X_{i+1}}+S_{i+1}+B_{i+1}+D_{i+1})> K_{X_i}+S_i+B_i+D_i$, which is not lc and we get a contradiction to ACC of lct.

    Therefore, we may assume that in the outcome $(X',S'+B'+D')$ of $-(K_X+B+S+D)$-MMP, neither $S$ nor $D$ is contracted. $(X', S'+B'+D')$ is not a Mori fiber space$/Z$ when $\epsilon$ is small enough by global ACC for for Fano type fibrations \cite[5.1]{waldron3mfsimperf}. So $-(K_{X'}+S'+B'+D')$ is nef$/Z$, we only need to deal with the outcome $(X',S'+D'+B')$. Finally, consider the semi-ample fibration $(X',S'+D'+B')\to V/Z$, if $V$ is a curve, then $-(K_{X'}+S'+D'+B')$ is trivial over $Z$ as $V\simeq Z$, and this is done by \Cref{relacomp1}; if $V$ is a surface, then this is done by \Cref{relacomp2}; if $X\to V$ is birational, then $-(K_{X'}+S'+D'+B')$ is nef and big$/Z$, and this is done by \Cref{relacomp4}.
\end{proof}
Using the same method as in \Cref{relacomp1} to construct a vertical lc center $S$, together with the results in \Cref{relacomp3}, \Cref{relacomp4} and \Cref{relacomp5}, we obtain the following result:
\begin{theorem}
\label{relacomp6}
    Assume $(X,B)$ is a 3-dimensional projective lc pair, $f:X\to Z$ is a contraction with $\dim Z=1$, let $R\subset [0,1]$ be a finite set of rational numbers, then there is some natural number $n=n(R)\in\mathbb N$ and a prime number $p_0=p_0(R)$ such that suppose:
    \begin{enumerate}
        \item $\chara k>p_0$,
        \item $B\in\Phi(R)$,
        \item $X$ is of Fano type$/Z$,
        \item $-(K_X+B)$ is nef and big$/Z$.
    \end{enumerate}
    Then, for any $z\in Z$, there is an $n$-complement $K_X+B^+$ of $K_X+B$ over $z$ with $B^+\geq B$.
\end{theorem}
Combining \Cref{relacomp1}, \Cref{relacomp2} and \Cref{relacomp6}, we finally obtain the boundedness of relative complements for 3-dimensional Fano type fibrations in large characteristic:
\begin{theorem}\label{recomp}
    Assume $(X,B)$ is a 3-dimensional projective lc pair, $f:X\to Z$ is a contraction with $3>\dim Z>0$, let $R\subset [0,1]$ be a finite set of rational numbers, then there is some natural number $n=n(R)\in\mathbb N$ a prime number $p_0=p_0(R)$ such that suppose:
    \begin{enumerate}
        \item $\chara k>p_0$,
        \item $B\in\Phi(R)$,
        \item $X$ is of Fano type$/Z$,
        \item $-(K_X+B)$ is nef$/Z$.
    \end{enumerate}
    Then, for any $z\in Z$, there is an $n$-complement $K_X+B^+$ of $K_X+B$ over $z$ with $B^+\geq B$.
\end{theorem}

\bibliographystyle{alpha}
\bibliography{cite}
\end{document}
